\chapter{Problems Related to Volume V}
\label{ch:companion-v}

\paragraph{Main-text correspondence.}
This Part accompanies \emph{Volume V: Commutative Algebra and Homological Methods},
whose canonical chapter range is \texttt{V/01--V/28}. Every distinct formal source
problem or exercise identified in the reconciled commutative-algebra corpus is preserved
as a canonical Companion problem; true duplicates and explanatory prose are not counted
as separate problems.

\paragraph{Organization.}
Problems are ordered by the exact main-volume chapter to which they most naturally
correspond. Fresh canonical problems are added only where the source corpus leaves a
chapter with zero or one dedicated problem.


\section*{V/01 --- Rings, Ideals and Quotients}

\begin{problem}[CP-V-0005]
\label{prob:cp-v-0005}
Use quotient rings to determine the indicated ideal-theoretic properties.

\textbf{(a)} Show that
\[
\mathbb Z[x]/(2,x)\cong\mathbb F_2
\]
and deduce that \((2,x)\) is maximal in \(\mathbb Z[x]\).

\textbf{(b)} Show that
\[
k[x,y]/(y-x^2)\cong k[x]
\]
and deduce that \((y-x^2)\) is prime.

\textbf{(c)} Show that \((xy)\subseteq k[x,y]\) is not prime.

\textbf{(d)} Show that \(k[x]/(x^2)\) is not reduced and compute
\(\sqrt{(x^2)}\).
\end{problem}

\begin{solution}
For \textbf{(a)}, evaluation at \(x=0\) followed by reduction modulo \(2\)
gives a surjective homomorphism
\[
\mathbb Z[x]\longrightarrow\mathbb F_2
\]
whose kernel is \((2,x)\). The first isomorphism theorem gives the claimed
quotient isomorphism. Since the quotient is a field, \((2,x)\) is maximal.

For \textbf{(b)}, define
\[
\psi:k[x,y]\longrightarrow k[x],
\qquad
\psi(f(x,y))=f(x,x^2).
\]
This map is surjective and its kernel is \((y-x^2)\). Hence
\[
k[x,y]/(y-x^2)\cong k[x].
\]
Because \(k[x]\) is a domain, the kernel is prime.

For \textbf{(c)}, the classes of \(x\) and \(y\) in
\(k[x,y]/(xy)\) are nonzero, but their product is zero. Thus the quotient
is not a domain, so \((xy)\) is not prime.

For \textbf{(d)}, the class \(\bar x\) is nonzero and satisfies
\(\bar x^2=0\), so the quotient is not reduced. Moreover
\[
f(x)^m\in(x^2)
\]
for some \(m\ge1\) exactly when \(f(0)=0\), which is equivalent to
\(f\in(x)\). Therefore
\[
\sqrt{(x^2)}=(x).
\]
\end{solution}


\begin{problem}[CP-V-0010]
\label{prob:cp-v-0010}
Let \(A\) be a ring, \(I\subseteq A\) an ideal, and let
\[
I[x]=\left\{a_0+a_1x+\cdots+a_nx^n\in A[x]:a_i\in I\text{ for all }i\right\}.
\]

\textbf{(a)} Prove that \(I[x]\) is exactly the extended ideal
\[
IA[x].
\]

\textbf{(b)} Construct a natural isomorphism
\[
(A/I)[x]\cong A[x]/I[x].
\]
\end{problem}

\begin{solution}
Every element of \(IA[x]\) is a finite sum of products \(if(x)\) with
\(i\in I\) and \(f(x)\in A[x]\), so all of its coefficients belong to \(I\).
Thus
\[
IA[x]\subseteq I[x].
\]
Conversely, if
\[
f(x)=a_0+a_1x+\cdots+a_nx^n
\]
has every \(a_i\in I\), then
\[
f(x)=a_0+a_1x+\cdots+a_nx^n
\]
is a sum of elements \(a_ix^i\in IA[x]\). Hence
\[
I[x]=IA[x].
\]

Now define
\[
\Phi:A[x]\longrightarrow (A/I)[x]
\]
by reducing every coefficient modulo \(I\):
\[
\Phi\!\left(\sum_i a_ix^i\right)=\sum_i(a_i+I)x^i.
\]
This is a surjective ring homomorphism. Its kernel consists exactly of
polynomials all of whose coefficients lie in \(I\), namely \(I[x]\).
The first isomorphism theorem therefore gives
\[
A[x]/I[x]\cong(A/I)[x].
\]
\end{solution}



\section*{V/02 --- Prime and Maximal Ideals}

\begin{problem}[CP-V-0002]
\label{prob:cp-v-0002}
Let
\[
\operatorname{Spec}A=\{\mathfrak p\subseteq A:\mathfrak p\text{ is prime}\},
\qquad
V(I)=\{\mathfrak p\in\operatorname{Spec}A:I\subseteq\mathfrak p\}.
\]

\textbf{(a)} Prove that the sets \(V(I)\) are precisely the closed sets of a topology
on \(\operatorname{Spec}A\).

\textbf{(b)} For \(f\in A\), put
\[
D(f)=\{\mathfrak p\in\operatorname{Spec}A:f\notin\mathfrak p\}.
\]
Prove that the sets \(D(f)\) form a basis.

\textbf{(c)} If \(\varphi:A\to B\) is a ring homomorphism, prove that
\[
\varphi^*:\operatorname{Spec}B\longrightarrow\operatorname{Spec}A,
\qquad
\mathfrak q\longmapsto\varphi^{-1}(\mathfrak q),
\]
is well defined and continuous.
\end{problem}

\begin{solution}
We have
\[
V(0)=\operatorname{Spec}A,\qquad V(A)=\varnothing.
\]
For any family \((I_\lambda)\),
\[
\bigcap_\lambda V(I_\lambda)=V\!\left(\sum_\lambda I_\lambda\right),
\]
while
\[
V(I)\cup V(J)=V(IJ)=V(I\cap J).
\]
Indeed, a prime ideal contains \(IJ\) exactly when it contains \(I\) or \(J\).
Thus the sets \(V(I)\) satisfy the closed-set axioms.

Since \(D(f)=\operatorname{Spec}A\setminus V((f))\), each \(D(f)\) is open.
Moreover,
\[
D(f)\cap D(g)=D(fg).
\]
If \(U=\operatorname{Spec}A\setminus V(I)\) and \(\mathfrak p\in U\), then
\(I\nsubseteq\mathfrak p\), so some \(f\in I\setminus\mathfrak p\). Hence
\[
\mathfrak p\in D(f)\subseteq U,
\]
which proves the basis assertion.

Finally, if \(\mathfrak q\in\operatorname{Spec}B\) and
\(ab\in\varphi^{-1}(\mathfrak q)\), then
\(\varphi(a)\varphi(b)\in\mathfrak q\); primality gives
\(a\in\varphi^{-1}(\mathfrak q)\) or \(b\in\varphi^{-1}(\mathfrak q)\).
Thus contraction preserves prime ideals. For an ideal \(I\subseteq A\),
\[
(\varphi^*)^{-1}(V(I))
=
V(IB),
\]
equivalently the preimage of a closed set is closed. Therefore \(\varphi^*\)
is continuous.
\end{solution}


\begin{problem}[CP-V-0006]
\label{prob:cp-v-0006}
Classify the maximal ideals in the following rings.

\textbf{(a)} \(A=\mathbb Z\).

\textbf{(b)} \(A=k[x]\), where \(k\) is a field.

\textbf{(c)} Specialize \textbf{(b)} when \(k\) is algebraically closed.
\end{problem}

\begin{solution}
Every ideal of \(\mathbb Z\) is principal. A proper ideal \((n)\) is maximal
exactly when
\[
\mathbb Z/(n)
\]
is a field, which occurs exactly when \(|n|=p\) is prime. Thus the maximal
ideals are \((p)\) for prime integers \(p\).

Since \(k[x]\) is a PID, every nonzero prime ideal is generated by an
irreducible polynomial. The quotient
\[
k[x]/(f)
\]
is a field exactly when \(f\) is irreducible, so the maximal ideals are
precisely \((f)\) with \(f\) irreducible.

If \(k\) is algebraically closed, every irreducible polynomial of positive
degree is linear. Hence the maximal ideals are exactly
\[
(x-a),\qquad a\in k.
\]
\end{solution}


\begin{problem}[CP-V-0007]
\label{prob:cp-v-0007}
Let \(k\) be algebraically closed. Prove that every maximal ideal of
\(k[x,y]\) has the form
\[
(x-a,y-b)
\]
for a unique point \((a,b)\in k^2\). Conversely, prove that every ideal of
this form is maximal.
\end{problem}

\begin{solution}
Let \(\mathfrak m\) be maximal. By the weak Nullstellensatz, the residue
field
\[
k[x,y]/\mathfrak m
\]
is a finite algebraic extension of \(k\). Since \(k\) is algebraically
closed, this extension is \(k\) itself.

Let \(a,b\in k\) be the images of \(x,y\) under the quotient map. Then
\(x-a\) and \(y-b\) belong to \(\mathfrak m\), so
\[
(x-a,y-b)\subseteq\mathfrak m.
\]
But evaluation at \((a,b)\) gives
\[
k[x,y]/(x-a,y-b)\cong k,
\]
hence \((x-a,y-b)\) is maximal. The inclusion between two maximal ideals
therefore forces equality.

Conversely, the same quotient isomorphism shows directly that every
\((x-a,y-b)\) is maximal.

Uniqueness follows because equality
\[
(x-a,y-b)=(x-a',y-b')
\]
implies both \(a-a'\) and \(b-b'\) vanish in the quotient, hence
\(a=a'\) and \(b=b'\).
\end{solution}


\begin{problem}[CP-V-0011]
\label{prob:cp-v-0011}
Let \(A\) be a ring and let \(\mathfrak p\subseteq A\) be a prime ideal.

\textbf{(a)} Prove that
\[
\mathfrak p[x]=\mathfrak p A[x]
\]
is a prime ideal of \(A[x]\).

\textbf{(b)} More generally, if \(I\subseteq A\), prove that
\[
A[x]/I[x]\cong(A/I)[x],
\]
and use this to show that \(I[x]\) is prime if and only if \(I\) is prime.

\textbf{(c)} Decide whether maximality is preserved: must
\(\mathfrak m[x]\) be maximal in \(A[x]\) when \(\mathfrak m\) is maximal
in \(A\)?
\end{problem}

\begin{solution}
By the coefficient-reduction isomorphism,
\[
A[x]/\mathfrak p[x]\cong(A/\mathfrak p)[x].
\]
Since \(\mathfrak p\) is prime, \(A/\mathfrak p\) is an integral domain.
A polynomial ring over a domain is again a domain, so the quotient
\(A[x]/\mathfrak p[x]\) is a domain. Therefore \(\mathfrak p[x]\) is prime.

The same argument for an arbitrary ideal \(I\) gives
\[
A[x]/I[x]\cong(A/I)[x].
\]
Now \((A/I)[x]\) is a domain exactly when \(A/I\) is a domain. Hence
\[
I[x]\text{ is prime}\quad\Longleftrightarrow\quad I\text{ is prime}.
\]

Maximality is not preserved in general. If \(\mathfrak m\) is maximal, then
\[
A[x]/\mathfrak m[x]\cong(A/\mathfrak m)[x].
\]
The coefficient ring \(A/\mathfrak m\) is a field, but its polynomial ring
in one indeterminate is not a field. Consequently \(\mathfrak m[x]\) is
prime but not maximal.
\end{solution}


\begin{problem}[CP-V-0090]
\label{prob:cp-v-0090}
Describe all points of \(\operatorname{Spec}(k[t])\).
\end{problem}

\begin{solution}
The prime ideals of \(k[t]\) are \((0)\) and the ideals \((f)\) generated by irreducible polynomials \(f\in k[t]\). Thus
\[
\operatorname{Spec}(k[t])=\{(0)\}\cup\{(f):f\text{ irreducible}\}.
\]
If \(k\) is algebraically closed, the nonzero prime ideals are exactly \((t-a)\), \(a\in k\). The generic point is \((0)\).
\end{solution}


\begin{problem}[CP-V-0091]
\label{prob:cp-v-0091}
Compute the residue fields of \((0)\), \((t-a)\), and \((t^2+1)\subset \mathbb R[t]\).
\end{problem}

\begin{solution}
For a prime \(\mathfrak p\), the residue field is
\[
\kappa(\mathfrak p)=\operatorname{Frac}(A/\mathfrak p).
\]
Hence
\[
\kappa((0))=\mathbb R(t),\qquad
\kappa((t-a))=\mathbb R,
\]
and
\[
\kappa((t^2+1))\cong \mathbb R[t]/(t^2+1)\cong\mathbb C.
\]
\end{solution}


\begin{problem}[CP-V-0093]
\label{prob:cp-v-0093}
Describe the closed sets of \(\operatorname{Spec}(\mathbb Z)\).
\end{problem}

\begin{solution}
The prime ideals are \((0)\) and \((p)\) for primes \(p\). For a nonzero integer \(n\),
\[
V((n))=\{(p):p\mid n\},
\]
a finite set of closed points. Also
\[
V((0))=\operatorname{Spec}(\mathbb Z),\qquad V((1))=\varnothing.
\]
Thus the proper closed subsets are precisely finite sets of closed points.
\end{solution}


\begin{problem}[CP-V-0094]
\label{prob:cp-v-0094}
Show that \(V(xy)=V(x)\cup V(y)\) in \(k[x,y]\).
\end{problem}

\begin{solution}
A prime ideal \(\mathfrak p\) contains \(xy\) if and only if, by primality, it contains \(x\) or \(y\). Therefore
\[
\mathfrak p\in V(xy)
\iff
\mathfrak p\in V(x)\cup V(y).
\]
\end{solution}



\section*{V/03 --- Radicals and Nilpotents}

\begin{problem}[CP-V-0004]
\label{prob:cp-v-0004}
Let
\[
A=k[x_1,\dots,x_n],
\]
where \(k\) is algebraically closed. For an ideal \(I\subseteq A\), define
\[
Z(I)=\{a\in k^n:f(a)=0\text{ for every }f\in I\},
\]
and for \(X\subseteq k^n\), define
\[
I(X)=\{f\in A:f(a)=0\text{ for every }a\in X\}.
\]

\textbf{(a)} Prove that every proper ideal \(I\ne A\) has \(Z(I)\ne\varnothing\).

\textbf{(b)} Prove the strong Nullstellensatz:
\[
I(Z(I))=\sqrt I.
\]
\end{problem}

\begin{solution}
For \textbf{(a)}, choose a maximal ideal \(\mathfrak m\supseteq I\).
By the weak Nullstellensatz,
\[
\mathfrak m=(x_1-a_1,\dots,x_n-a_n)
\]
for some \(a=(a_1,\dots,a_n)\in k^n\). Every element of \(I\) vanishes at
\(a\), so \(a\in Z(I)\).

For \textbf{(b)}, if \(f\in\sqrt I\), then \(f^m\in I\) for some \(m\).
Thus \(f(a)^m=0\) for every \(a\in Z(I)\), and hence \(f(a)=0\). Therefore
\[
\sqrt I\subseteq I(Z(I)).
\]

Conversely, let \(f\in I(Z(I))\). In
\(k[x_1,\dots,x_n,y]\), set
\[
J=(I,1-yf).
\]
The set \(Z(J)\) is empty: a common zero would satisfy both \(f(a)=0\) and
\(1-yf(a)=0\). Part \textbf{(a)} therefore implies \(J=(1)\). Hence
\[
1=g+h(1-yf)
\]
with \(g\in I\,k[x_1,\dots,x_n,y]\). Substituting \(y=f^{-1}\) in the
fraction field and clearing a sufficiently large power of \(f\) gives
\[
f^N\in I
\]
for some \(N\ge1\). Thus \(f\in\sqrt I\), proving
\[
I(Z(I))=\sqrt I.
\]
\end{solution}


\begin{problem}[CP-V-0092]
\label{prob:cp-v-0092}
Compare \(k[x]/(x)\) and \(k[x]/(x^2)\). Do they have the same topological space? Do they define the same scheme?
\end{problem}

\begin{solution}
Both spectra have one prime ideal, so their underlying topological spaces are singletons. But
\[
k[x]/(x)\cong k
\]
is reduced, whereas \(k[x]/(x^2)\) contains the nonzero nilpotent class of \(x\). Thus the schemes are not isomorphic even though their underlying spaces agree.
\end{solution}


\begin{problem}[CP-V-0101]
\label{prob:cp-v-0101}
In \(A=k[x]/(x^2)\), compute \(D(\bar x)\).
\end{problem}

\begin{solution}
Every prime ideal of \(A\) contains the nilpotent element \(\bar x\), since \(\bar x^2=0\). Therefore no prime avoids \(\bar x\), and
\[
D(\bar x)=\varnothing.
\]
\end{solution}



\section*{V/04 --- Chinese Remainder Theory}

\begin{problem}[CP-V-0001]
\label{prob:cp-v-0001}
Let \(A\) be a ring and let \(I_1,\dots,I_n\) be ideals. Define
\[
\phi:A\longrightarrow \prod_{i=1}^n A/I_i,
\qquad
\phi(a)=(a+I_1,\dots,a+I_n).
\]
Two ideals \(I,J\) are called coprime if \(I+J=A\).

\textbf{(a)} Assume \(I_i+I_j=A\) whenever \(i\ne j\). Prove that
\[
\bigcap_{i=1}^n I_i=\prod_{i=1}^n I_i.
\]

\textbf{(b)} Prove that \(\phi\) is surjective if and only if the ideals are pairwise coprime.

\textbf{(c)} Prove that \(\phi\) is injective if and only if
\[
\bigcap_{i=1}^n I_i=0.
\]
\end{problem}

\begin{solution}
For two coprime ideals \(I,J\), choose \(x\in I\) and \(y\in J\) with \(x+y=1\).
The inclusion \(IJ\subseteq I\cap J\) is automatic. If \(z\in I\cap J\), then
\[
z=z(x+y)=zx+zy\in IJ,
\]
so \(I\cap J=IJ\).

\textbf{(a)} Proceed by induction. Suppose
\[
\bigcap_{i=1}^k I_i=\prod_{i=1}^k I_i=:P_k.
\]
For each \(i\le k\), choose \(a_i\in I_i\) and \(b_i\in I_{k+1}\) such that
\(a_i+b_i=1\). Expanding \(\prod_{i=1}^k(a_i+b_i)\) shows that
\(1=u+v\) with \(u\in P_k\) and \(v\in I_{k+1}\). Hence
\(P_k+I_{k+1}=A\), and therefore
\[
\bigcap_{i=1}^{k+1}I_i
=P_k\cap I_{k+1}
=P_kI_{k+1}
=\prod_{i=1}^{k+1}I_i.
\]

\textbf{(b)} If \(\phi\) is surjective, then its projection to
\(A/I_i\times A/I_j\) is surjective for every \(i\ne j\). The two-ideal
Chinese remainder theorem gives \(I_i+I_j=A\).

Conversely, assume pairwise coprimality. For fixed \(i\), put
\[
J_i=\bigcap_{j\ne i} I_j.
\]
The same product argument gives \(I_i+J_i=A\), so choose \(e_i\in J_i\)
with \(e_i\equiv1\pmod{I_i}\). Then \(e_i\equiv0\pmod{I_j}\) for
\(j\ne i\). Given residue classes \(\overline a_i\in A/I_i\), the element
\[
a=\sum_{i=1}^n a_i e_i
\]
maps to \((\overline a_1,\dots,\overline a_n)\). Thus \(\phi\) is surjective.

\textbf{(c)} Directly,
\[
\ker\phi
=\{a\in A:a\in I_i\text{ for every }i\}
=\bigcap_{i=1}^n I_i.
\]
Hence \(\phi\) is injective exactly when the intersection is zero.
\end{solution}


\begin{problem}[CP-V-0003]
\label{prob:cp-v-0003}
Let \(A\) be a ring. Prove that the following conditions are equivalent.

\textbf{(i)} \(\operatorname{Spec}A\) is disconnected.

\textbf{(ii)} There exist nonzero elements \(e_1,e_2\in A\) such that
\[
e_1^2=e_1,\qquad e_2^2=e_2,\qquad e_1e_2=0,\qquad e_1+e_2=1.
\]

\textbf{(iii)} There exist nonzero rings \(A_1,A_2\) with
\[
A\cong A_1\times A_2.
\]
\end{problem}

\begin{solution}
Assume \textbf{(ii)} and set \(A_i=Ae_i\). The map
\[
A\longrightarrow A_1\times A_2,
\qquad
a\longmapsto(ae_1,ae_2),
\]
is injective because \(a=a(e_1+e_2)\), and it is surjective because
\((ae_1,be_2)\) is the image of \(ae_1+be_2\). Hence \textbf{(iii)} holds.

If \(A\cong A_1\times A_2\) with both factors nonzero, then
\[
\operatorname{Spec}A
\cong
\operatorname{Spec}A_1\sqcup\operatorname{Spec}A_2.
\]
Both pieces are nonempty, open, and closed, so \(\operatorname{Spec}A\) is
disconnected. Thus \textbf{(iii)} implies \textbf{(i)}.

Now suppose
\[
\operatorname{Spec}A=U\sqcup V
\]
with \(U,V\) nonempty and both open and closed. Write
\[
U=V(I),\qquad V=V(J)
\]
for ideals \(I,J\). Since \(U\cap V=\varnothing\),
\[
V(I+J)=\varnothing,
\]
so \(I+J=A\). Since \(U\cup V=\operatorname{Spec}A\),
\[
V(IJ)=\operatorname{Spec}A,
\]
hence \(IJ\subseteq\sqrt{(0)}\). Passing first to the reduced quotient
\(A_{\mathrm{red}}=A/\sqrt{(0)}\), the images of \(I\) and \(J\) are
comaximal with zero product. Chinese remainder theory then gives
\[
A_{\mathrm{red}}\cong A_{\mathrm{red}}/\overline I
\times A_{\mathrm{red}}/\overline J,
\]
and therefore a nontrivial idempotent \(\overline e\in A_{\mathrm{red}}\).
Idempotents lift uniquely across nilpotent ideals, so \(\overline e\) lifts
to an idempotent \(e\in A\) with \(e\ne0,1\). Taking
\[
e_1=e,\qquad e_2=1-e
\]
gives \textbf{(ii)}.
\end{solution}


\begin{problem}[CP-V-0008]
\label{prob:cp-v-0008}
Let \(k\) be a field and let \(a,b,c\in k\) be pairwise distinct. Prove that
the ideals
\[
(x-a),\qquad(x-b),\qquad(x-c)
\]
of \(k[x]\) are pairwise comaximal and deduce an isomorphism
\[
k[x]/\bigl((x-a)(x-b)(x-c)\bigr)
\cong
k\times k\times k.
\]
Describe the isomorphism explicitly.
\end{problem}

\begin{solution}
For \(a\ne b\),
\[
(x-a)-(x-b)=b-a.
\]
Since \(b-a\in k^\times\), the ideal generated by \(x-a\) and \(x-b\)
contains a unit. Hence
\[
(x-a)+(x-b)=k[x].
\]
The same argument works for every pair.

The Chinese remainder theorem therefore gives
\[
k[x]\Big/\bigl((x-a)\cap(x-b)\cap(x-c)\bigr)
\cong
k[x]/(x-a)\times k[x]/(x-b)\times k[x]/(x-c).
\]
Pairwise comaximality also yields
\[
(x-a)\cap(x-b)\cap(x-c)
=
\bigl((x-a)(x-b)(x-c)\bigr).
\]
Each quotient by a linear factor is naturally \(k\), so the desired
isomorphism is
\[
[f]\longmapsto\bigl(f(a),f(b),f(c)\bigr).
\]
\end{solution}


\begin{problem}[CP-V-0009]
\label{prob:cp-v-0009}
Compute \(I\cap J\) and \(IJ\) in each case, and decide whether
\(I\cap J=IJ\).

\textbf{(a)} \(I=(2)\), \(J=(3)\) in \(\mathbb Z\).

\textbf{(b)} \(I=(4)\), \(J=(6)\) in \(\mathbb Z\).

\textbf{(c)} \(I=(x)\), \(J=(x-1)\) in \(k[x]\).

\textbf{(d)} \(I=(x)\), \(J=(x^2)\) in \(k[x]\).

Explain the relation with comaximality.
\end{problem}

\begin{solution}
For principal ideals in \(\mathbb Z\),
\[
(m)\cap(n)=(\operatorname{lcm}(m,n)),
\qquad
(m)(n)=(mn).
\]
Thus
\[
(2)\cap(3)=(6)=(2)(3),
\]
whereas
\[
(4)\cap(6)=(12)\ne(24)=(4)(6).
\]

In \(k[x]\), the ideals \((x)\) and \((x-1)\) are comaximal because
\[
x-(x-1)=1.
\]
Hence
\[
(x)\cap(x-1)=(x(x-1))=(x)(x-1).
\]

Since \((x^2)\subseteq(x)\),
\[
(x)\cap(x^2)=(x^2),
\]
whereas
\[
(x)(x^2)=(x^3).
\]
So equality fails.

In general, comaximal ideals \(I+J=A\) satisfy
\[
I\cap J=IJ.
\]
The noncomaximal examples above show that this equality need not hold
without the hypothesis.
\end{solution}



\section*{V/05 --- Multiplicative Systems}

\begin{problem}[CP-V-0102]
\label{prob:cp-v-0102}
Prove that \(D(fg)=D(f)\cap D(g)\).
\end{problem}

\begin{solution}
For a prime \(\mathfrak p\),
\[
fg\notin\mathfrak p
\iff
f\notin\mathfrak p\ \text{and}\ g\notin\mathfrak p,
\]
because if either factor belongs to \(\mathfrak p\), so does the product, and if the product belongs to a prime, at least one factor does. Hence
\[
D(fg)=D(f)\cap D(g).
\]
\end{solution}


\begin{problem}[CP-V-0103]
\label{prob:cp-v-0103}
Prove that \(D(1)=\operatorname{Spec}(A)\) and \(D(0)=\varnothing\).
\end{problem}

\begin{solution}
No proper prime ideal contains \(1\), so every prime lies in \(D(1)\). Every prime contains \(0\), so none lies in \(D(0)\). Thus
\[
D(1)=\operatorname{Spec}(A),\qquad D(0)=\varnothing.
\]
\end{solution}


\begin{problem}[CP-V-0104]
\label{prob:cp-v-0104}
Show that \(D(f)=D(f^2)\).
\end{problem}

\begin{solution}
For a prime ideal \(\mathfrak p\),
\[
f^2\in\mathfrak p\iff f\in\mathfrak p.
\]
Taking complements gives
\[
D(f)=D(f^2).
\]
\end{solution}



\section*{V/06 --- Localization of Rings}

\begin{problem}[CP-V-0013]
\label{prob:cp-v-0013}
Let
\[
R=\frac{k[x,y,z]}{(xy-z^2)}.
\]
Prove that localization at the powers of \(x\) gives
\[
R_x\cong k[x,z]_x.
\]
Explain algebraically why the variable \(y\) becomes redundant after inverting
\(x\).
\end{problem}

\begin{solution}
In \(R_x\), the class of \(x\) is invertible. Since the defining relation is
\[
xy=z^2,
\]
we obtain
\[
y=\frac{z^2}{x}.
\]
Thus \(R_x\) is generated by \(x,x^{-1},z\).

Define
\[
\varphi:k[x,z]_x\longrightarrow R_x
\]
by
\[
\varphi(x)=\bar x,\qquad \varphi(z)=\bar z.
\]
This is surjective because
\[
\bar y=\frac{\bar z^2}{\bar x}
\]
lies in its image.

Conversely, define
\[
\psi:R\longrightarrow k[x,z]_x
\]
by
\[
\psi(x)=x,\qquad
\psi(z)=z,\qquad
\psi(y)=\frac{z^2}{x}.
\]
The relation is respected:
\[
\psi(xy-z^2)
=
x\frac{z^2}{x}-z^2
=
0.
\]
Hence \(\psi\) factors through \(R\). Since \(x\) maps to a unit, the universal
property of localization extends it to
\[
\psi_x:R_x\longrightarrow k[x,z]_x.
\]
The two maps are inverse, so
\[
R_x\cong k[x,z]_x.
\]
The redundancy of \(y\) is exactly the fact that after \(x\) becomes a unit,
the relation solves uniquely for \(y\).
\end{solution}


\begin{problem}[CP-V-0014]
\label{prob:cp-v-0014}
Let \(\varphi:A\to B\) be a ring homomorphism, with induced map
\[
\varphi^*:\operatorname{Spec}B\to\operatorname{Spec}A.
\]

\textbf{(a)} Let \(S\subseteq A\) be multiplicatively closed. Prove that
\[
\operatorname{Spec}(S^{-1}A)
\]
is naturally homeomorphic to
\[
\{\mathfrak p\in\operatorname{Spec}A:\mathfrak p\cap S=\varnothing\}.
\]
In the principal case \(S=\{1,f,f^2,\dots\}\), identify this subspace with
\(D(f)\).

\textbf{(b)} Show that localization is compatible with \(\varphi^*\): if
\(S^{-1}X\subseteq\operatorname{Spec}A\) denotes the subspace above, then
\[
(\varphi^*)^{-1}(S^{-1}X)
\]
is the corresponding localization subspace of \(\operatorname{Spec}B\).

\textbf{(c)} Let \(I\subseteq A\) and let \(J=IB\). Prove that the map
\[
\operatorname{Spec}(B/J)\longrightarrow\operatorname{Spec}(A/I)
\]
induced by \(A/I\to B/J\) agrees with the restriction of \(\varphi^*\) to
\(V(J)\).

\textbf{(d)} Let \(\mathfrak p\in\operatorname{Spec}A\). Prove that the fiber of
\(\varphi^*\) over \(\mathfrak p\) is naturally
\[
\operatorname{Spec}
\left(
B_{\mathfrak p}/\mathfrak pB_{\mathfrak p}
\right),
\]
and hence
\[
(\varphi^*)^{-1}(\mathfrak p)
\cong
\operatorname{Spec}\bigl(k(\mathfrak p)\otimes_A B\bigr),
\]
where
\[
k(\mathfrak p)=A_{\mathfrak p}/\mathfrak pA_{\mathfrak p}.
\]
\end{problem}

\begin{solution}
For \textbf{(a)}, contraction gives a map
\[
\operatorname{Spec}(S^{-1}A)\longrightarrow\operatorname{Spec}A.
\]
If \(\mathfrak P\subseteq S^{-1}A\) is prime, then its contraction
\(\mathfrak p\) is disjoint from \(S\), because every element of \(S\) becomes
a unit after localization.

Conversely, if \(\mathfrak p\cap S=\varnothing\), then
\[
S^{-1}\mathfrak p
\]
is prime in \(S^{-1}A\), and extension and contraction are inverse. Thus there
is a bijection
\[
\operatorname{Spec}(S^{-1}A)
\cong
\{\mathfrak p\in\operatorname{Spec}A:\mathfrak p\cap S=\varnothing\}.
\]
A basic open
\[
D(a/s)\subseteq\operatorname{Spec}(S^{-1}A)
\]
corresponds to
\[
D(a)\cap\{\mathfrak p:\mathfrak p\cap S=\varnothing\},
\]
so the bijection is a homeomorphism onto the indicated subspace. If
\(S=\{1,f,f^2,\dots\}\), disjointness from \(S\) is equivalent to
\(f\notin\mathfrak p\), hence the image is \(D(f)\).

For \textbf{(b)}, the image \(\varphi(S)\) is multiplicatively closed in \(B\),
and the universal property gives
\[
S^{-1}A\longrightarrow\varphi(S)^{-1}B.
\]
For \(\mathfrak q\in\operatorname{Spec}B\),
\[
\varphi^{-1}(\mathfrak q)\cap S=\varnothing
\]
if and only if
\[
\mathfrak q\cap\varphi(S)=\varnothing.
\]
Therefore the inverse image of the localization subspace in
\(\operatorname{Spec}A\) is exactly the corresponding localization subspace in
\(\operatorname{Spec}B\).

For \textbf{(c)}, primes of \(A/I\) correspond to primes of \(A\) containing
\(I\), and primes of \(B/J\) correspond to primes of \(B\) containing \(J\).
If \(\mathfrak q\supseteq J\), then
\[
(\mathfrak q/J)\longmapsto \varphi^{-1}(\mathfrak q)/I.
\]
Under the standard identifications with \(V(J)\) and \(V(I)\), this is exactly
the restriction of \(\varphi^*\).

For \textbf{(d)}, localize at
\[
S=A\setminus\mathfrak p.
\]
Then
\[
S^{-1}A=A_{\mathfrak p},
\qquad
S^{-1}B=B_{\mathfrak p}.
\]
A prime \(\mathfrak q\subseteq B\) lies over \(\mathfrak p\) exactly when its
localized prime in \(B_{\mathfrak p}\) contains
\[
\mathfrak pB_{\mathfrak p}.
\]
Hence
\[
(\varphi^*)^{-1}(\mathfrak p)
\cong
\operatorname{Spec}
\left(
B_{\mathfrak p}/\mathfrak pB_{\mathfrak p}
\right).
\]
Finally,
\[
B_{\mathfrak p}/\mathfrak pB_{\mathfrak p}
\cong
k(\mathfrak p)\otimes_A B,
\]
so
\[
(\varphi^*)^{-1}(\mathfrak p)
\cong
\operatorname{Spec}\bigl(k(\mathfrak p)\otimes_A B\bigr).
\]
\end{solution}


\begin{problem}[CP-V-0015]
\label{prob:cp-v-0015}
Let
\[
A=k[x,y]/(xy),
\]
and write \(\bar x,\bar y\) for the residue classes of \(x,y\). Let
\[
S=\{1,\bar x,\bar x^2,\bar x^3,\dots\}.
\]

\textbf{(a)} Prove that
\[
S^{-1}A\cong k[x,x^{-1}].
\]

\textbf{(b)} Determine which of the prime ideals
\[
(\bar x),\qquad (\bar y),\qquad (\bar x,\bar y)
\]
survive under localization.

\textbf{(c)} Explain geometrically what localization at \(\bar x\) does to
the affine scheme defined by \(xy=0\).
\end{problem}

\begin{solution}
In \(A\) we have
\[
\bar x\bar y=0.
\]
After localizing at \(S\), the element \(\bar x\) becomes invertible, so
\[
\bar y=\bar x^{-1}(\bar x\bar y)=0.
\]
Hence every element of \(S^{-1}A\) is represented using only \(\bar x\) and
\(\bar x^{-1}\). This gives a natural surjective map
\[
k[x,x^{-1}]\longrightarrow S^{-1}A.
\]
Conversely, the homomorphism
\[
k[x,y]\longrightarrow k[x,x^{-1}],
\qquad
x\mapsto x,\qquad y\mapsto 0
\]
kills \(xy\), hence factors through \(A\), and since \(x\) maps to a unit it
extends uniquely to \(S^{-1}A\). The two maps are inverse, so
\[
S^{-1}A\cong k[x,x^{-1}].
\]

A prime \(\mathfrak p\subseteq A\) survives precisely when
\[
\mathfrak p\cap S=\varnothing.
\]
Thus \((\bar x)\) and \((\bar x,\bar y)\) disappear because they contain
\(\bar x\), while \((\bar y)\) survives.

Geometrically, \(xy=0\) is the union of the coordinate axes. Localizing at
\(\bar x\) restricts to the open set \(D(\bar x)\), so the component
\(x=0\) disappears, as does the intersection point. What remains is the
\(x\)-axis away from the origin.
\end{solution}


\begin{problem}[CP-V-0059]
\label{prob:cp-v-0059}
Determine whether
\[
\frac{1}{x}
\]
belongs to
\[
k[x]_x
\qquad\text{and}\qquad
k[x]_{(x)}.
\]
\end{problem}

\begin{solution}
In
\[
k[x]_x,
\]
the allowed denominators are powers of \(x\). Since
\[
x=x^1,
\]
we have
\[
\frac{1}{x}\in k[x]_x.
\]

In
\[
k[x]_{(x)},
\]
the allowed denominators are polynomials not in \((x)\), that is, polynomials with nonzero constant term. But
\[
x\in (x),
\]
so
\[
\frac{1}{x}\notin k[x]_{(x)}.
\]
\end{solution}


\begin{problem}[CP-V-0060]
\label{prob:cp-v-0060}
Determine whether
\[
\frac{1}{x+1}
\]
belongs to
\[
k[x]_x
\qquad\text{and}\qquad
k[x]_{(x)}.
\]
\end{problem}

\begin{solution}
In
\[
k[x]_x,
\]
the denominator must be a power of \(x\). Since \(x+1\) is not a power of \(x\),
\[
\frac{1}{x+1}\notin k[x]_x.
\]

In
\[
k[x]_{(x)},
\]
the denominator is allowed if it does not lie in \((x)\). Since
\[
(x+1)(0)=1\neq 0,
\]
we have
\[
x+1\notin (x),
\]
hence
\[
\frac{1}{x+1}\in k[x]_{(x)}.
\]
\end{solution}


\begin{problem}[CP-V-0061]
\label{prob:cp-v-0061}
Determine whether
\[
\frac{x^2+1}{x-2}
\]
belongs to
\[
k[x]_{(x)}
\qquad\text{and}\qquad
k[x]_{(x-2)}.
\]
\end{problem}

\begin{solution}
For
\[
k[x]_{(x)},
\]
the denominator must satisfy
\[
(x-2)(0)=-2\neq 0.
\]
So
\[
x-2\notin (x),
\]
hence
\[
\frac{x^2+1}{x-2}\in k[x]_{(x)}.
\]

For
\[
k[x]_{(x-2)},
\]
the denominator must not lie in \((x-2)\). But
\[
x-2\in (x-2),
\]
so
\[
\frac{x^2+1}{x-2}\notin k[x]_{(x-2)}.
\]
\end{solution}


\begin{problem}[CP-V-0062]
\label{prob:cp-v-0062}
Determine whether
\[
\frac{x^2+1}{x^2-1}
\]
belongs to
\[
k[x]_{(x-1)}.
\]
\end{problem}

\begin{solution}
The denominator is
\[
x^2-1=(x-1)(x+1).
\]
Since it vanishes at \(x=1\),
\[
x^2-1\in (x-1).
\]
Therefore
\[
\frac{x^2+1}{x^2-1}\notin k[x]_{(x-1)}.
\]
\end{solution}


\begin{problem}[CP-V-0063]
\label{prob:cp-v-0063}
Determine whether
\[
\frac{x^2+1}{x^2-1}
\]
belongs to
\[
k[x]_{(x)}.
\]
\end{problem}

\begin{solution}
We test the denominator at \(x=0\):
\[
x^2-1\big|_{x=0}=-1\neq 0.
\]
So
\[
x^2-1\notin (x),
\]
and hence
\[
\frac{x^2+1}{x^2-1}\in k[x]_{(x)}.
\]
\end{solution}


\begin{problem}[CP-V-0064]
\label{prob:cp-v-0064}
Determine whether
\[
\frac{x^3+x+1}{x^5-7x+2}
\]
belongs to
\[
k[x]_{(0)}.
\]
\end{problem}

\begin{solution}
Since
\[
k[x]_{(0)}=k(x),
\]
every fraction with nonzero polynomial denominator belongs to this localization. So if
\[
x^5-7x+2\neq 0
\]
as a polynomial, then
\[
\frac{x^3+x+1}{x^5-7x+2}\in k[x]_{(0)}.
\]
\end{solution}


\begin{problem}[CP-V-0065]
\label{prob:cp-v-0065}
Determine whether
\[
\frac{1}{x(x+1)}
\]
belongs to
\[
k[x]_x
\qquad\text{and}\qquad
k[x]_{(0)}.
\]
\end{problem}

\begin{solution}
In
\[
k[x]_x,
\]
the denominator must be a power of \(x\). But
\[
x(x+1)
\]
is not a power of \(x\), so
\[
\frac{1}{x(x+1)}\notin k[x]_x.
\]

In
\[
k[x]_{(0)}=k(x),
\]
every nonzero polynomial denominator is allowed, so
\[
\frac{1}{x(x+1)}\in k[x]_{(0)}.
\]
\end{solution}


\begin{problem}[CP-V-0069]
\label{prob:cp-v-0069}
Determine whether
\[
\frac{x}{y}
\]
belongs to
\[
k[x,y]_{(y)}
\qquad\text{and}\qquad
k[x,y]_y.
\]
\end{problem}

\begin{solution}
In
\[
k[x,y]_{(y)},
\]
we invert all elements not in \((y)\). But
\[
y\in (y),
\]
so
\[
\frac{x}{y}\notin k[x,y]_{(y)}.
\]

In
\[
k[x,y]_y,
\]
we explicitly invert powers of \(y\). Since the denominator is \(y\),
\[
\frac{x}{y}\in k[x,y]_y.
\]
\end{solution}


\begin{problem}[CP-V-0070]
\label{prob:cp-v-0070}
Determine whether
\[
\frac{1}{xy}
\]
belongs to
\[
k[x,y]_x,
\qquad
k[x,y]_y,
\qquad
k[x,y]_{(x,y)}.
\]
\end{problem}

\begin{solution}
In
\[
k[x,y]_x,
\]
only powers of \(x\) are allowed as denominators. Since
\[
xy
\]
is not a power of \(x\),
\[
\frac{1}{xy}\notin k[x,y]_x.
\]

Similarly,
\[
\frac{1}{xy}\notin k[x,y]_y,
\]
because \(xy\) is not a power of \(y\).

In
\[
k[x,y]_{(x,y)},
\]
the denominator must not vanish at \((0,0)\). But
\[
xy(0,0)=0,
\]
so
\[
\frac{1}{xy}\notin k[x,y]_{(x,y)}.
\]
\end{solution}


\begin{problem}[CP-V-0073]
\label{prob:cp-v-0073}
Is
\[
\frac{1}{x}
\]
an element of
\[
k[x]/(x)?
\]
\end{problem}

\begin{solution}
No. The ring
\[
k[x]/(x)
\]
is just
\[
k.
\]
In this quotient one has
\[
x=0,
\]
so the symbol
\[
\frac{1}{x}
\]
does not make sense there. Quotients do not create fractions; they only impose relations. Thus
\[
\frac{1}{x}\notin k[x]/(x).
\]
\end{solution}


\begin{problem}[CP-V-0074]
\label{prob:cp-v-0074}
Compare
\[
\frac{1}{x},
\qquad
\frac{1}{x+1}
\]
in
\[
k[x]_x,
\qquad
k[x]_{(x)},
\qquad
k[x]_{(0)}.
\]
\end{problem}

\begin{solution}
For
\[
\frac{1}{x}:
\]
\[
\frac{1}{x}\in k[x]_x,
\qquad
\frac{1}{x}\notin k[x]_{(x)},
\qquad
\frac{1}{x}\in k[x]_{(0)}.
\]

For
\[
\frac{1}{x+1}:
\]
\[
\frac{1}{x+1}\notin k[x]_x,
\qquad
\frac{1}{x+1}\in k[x]_{(x)},
\qquad
\frac{1}{x+1}\in k[x]_{(0)}.
\]

This exercise clearly shows the differences among the three localizations.
\end{solution}


\begin{problem}[CP-V-0097]
\label{prob:cp-v-0097}
Describe \(D(t)\subset \operatorname{Spec}(k[t])\) and compute its ring of sections.
\end{problem}

\begin{solution}
By definition
\[
D(t)=\{\mathfrak p:t\notin\mathfrak p\}.
\]
It is the affine open
\[
D(t)\cong\operatorname{Spec}(k[t]_t)
=\operatorname{Spec}(k[t,t^{-1}]).
\]
Hence
\[
\Gamma(D(t),\mathcal O)=k[t,t^{-1}].
\]
\end{solution}


\begin{problem}[CP-V-0098]
\label{prob:cp-v-0098}
Show that \(D(f)\cong \operatorname{Spec}(A_f)\).
\end{problem}

\begin{solution}
Prime ideals of \(A_f\) correspond bijectively to prime ideals \(\mathfrak p\subseteq A\) disjoint from \(\{1,f,f^2,\dots\}\), equivalently to primes with \(f\notin\mathfrak p\). These are exactly the points of \(D(f)\). The correspondence preserves the Zariski topology, giving
\[
D(f)\cong\operatorname{Spec}(A_f).
\]
\end{solution}



\section*{V/07 --- Localization of Modules}

\begin{problem}[CP-V-0012]
\label{prob:cp-v-0012}
Let \(A\) be a ring.

\textbf{(a)} Let \(S\subseteq A\) be multiplicatively closed and let \(M\) be a
finitely generated \(A\)-module. Prove that
\[
S^{-1}M=0
\]
if and only if there exists \(s\in S\) such that
\[
sM=0.
\]

\textbf{(b)} Let \(S,T\subseteq A\) be multiplicatively closed, and let \(U\) be
the image of \(T\) in \(S^{-1}A\). Prove that
\[
(ST)^{-1}A\cong U^{-1}(S^{-1}A),
\qquad
ST=\{st:s\in S,\ t\in T\}.
\]

\textbf{(c)} Let \(f:A\to B\) be a ring homomorphism, let \(S\subseteq A\) be
multiplicatively closed, and put \(T=f(S)\). Prove that
\[
S^{-1}B\cong T^{-1}B
\]
as \(S^{-1}A\)-modules.
\end{problem}

\begin{solution}
For \textbf{(a)}, if \(sM=0\) for some \(s\in S\), then for every \(m\in M\),
\[
\frac{m}{1}=\frac{sm}{s}=0,
\]
so \(S^{-1}M=0\).

Conversely, assume \(S^{-1}M=0\), and let \(m_1,\dots,m_r\) generate \(M\).
For each \(i\),
\[
\frac{m_i}{1}=0
\]
in \(S^{-1}M\), so there exists \(s_i\in S\) such that \(s_i m_i=0\).
Set
\[
s=s_1\cdots s_r\in S.
\]
Then \(sm_i=0\) for every generator, hence \(sM=0\).

For \textbf{(b)}, define
\[
\Phi:(ST)^{-1}A\longrightarrow U^{-1}(S^{-1}A),
\qquad
\Phi\!\left(\frac{a}{st}\right)=\frac{a/s}{t},
\]
and
\[
\Psi:U^{-1}(S^{-1}A)\longrightarrow(ST)^{-1}A,
\qquad
\Psi\!\left(\frac{a/s}{t}\right)=\frac{a}{st}.
\]
These are well-defined ring homomorphisms and are inverse to one another. Thus
\[
(ST)^{-1}A\cong U^{-1}(S^{-1}A).
\]

For \textbf{(c)}, define
\[
\alpha:S^{-1}B\longrightarrow T^{-1}B,
\qquad
\alpha\!\left(\frac{b}{s}\right)=\frac{b}{f(s)}.
\]
The natural map \(B\to S^{-1}B\) sends every \(f(s)\in T\) to a unit, so by
the universal property of localization it extends uniquely to
\[
\beta:T^{-1}B\longrightarrow S^{-1}B,
\qquad
\beta\!\left(\frac{b}{f(s)}\right)=\frac{b}{s}.
\]
The maps \(\alpha\) and \(\beta\) are inverse \(S^{-1}A\)-linear
isomorphisms.
\end{solution}


\begin{problem}[CP-V-0080]
\label{prob:cp-v-0080}
What is the localization of \(\mathbb Z/15\mathbb Z\) at \((2)\)?
\end{problem}

\begin{solution}
Start with
\[
0 \longrightarrow \mathbb Z \xrightarrow{\cdot 15} \mathbb Z \longrightarrow \mathbb Z/15\mathbb Z \longrightarrow 0.
\]
After localizing at \((2)\), the element \(15\) becomes invertible in \(\mathbb Z_{(2)}\), since \(2\nmid 15\).
Therefore multiplication by \(15\) is an isomorphism, and the quotient becomes zero:
\[
(\mathbb Z/15\mathbb Z)_{(2)}=0.
\]
\end{solution}


\begin{problem}[CP-V-0081]
\label{prob:cp-v-0081}
What is the localization of \(\mathbb Z/15\mathbb Z\) at \((3)\)?
\end{problem}

\begin{solution}
In \(\mathbb Z_{(3)}\), the element \(5\) is invertible, so multiplication by \(15=3\cdot 5\) is equivalent, up to a unit, to multiplication by \(3\). Therefore
\[
(\mathbb Z/15\mathbb Z)_{(3)} \cong \mathbb Z_{(3)}/(3)\cong \mathbb F_3.
\]
\end{solution}



\section*{V/08 --- Local Rings and Localization at Primes}

\begin{problem}[CP-V-0016]
\label{prob:cp-v-0016}
Let \(\mathfrak p\) be a prime ideal of a ring \(A\), and set
\[
S=A\setminus\mathfrak p.
\]

\textbf{(a)} Prove that
\[
A_{\mathfrak p}=S^{-1}A
\]
is a local ring.

\textbf{(b)} Prove that its unique maximal ideal is
\[
\mathfrak pA_{\mathfrak p}.
\]

\textbf{(c)} Specialize to
\[
A=k[x,y],\qquad \mathfrak p=(x,y),
\]
and characterize the units of \(A_{\mathfrak p}\).

\textbf{(d)} Specialize to
\[
A=\mathbb Z,\qquad \mathfrak p=(p),
\]
where \(p\) is prime, and show that
\[
\mathbb Z_{(p)}
=
\left\{\frac ab\in\mathbb Q:p\nmid b\right\}
\]
has maximal ideal \(p\mathbb Z_{(p)}\).
\end{problem}

\begin{solution}
An element of \(A_{\mathfrak p}\) has the form \(a/s\) with
\(s\notin\mathfrak p\). If \(a\notin\mathfrak p\), then \(a\in S\), so
\[
\frac{a}{s}
\]
is a unit with inverse \(s/a\).

If \(a\in\mathfrak p\), then \(a/s\in\mathfrak pA_{\mathfrak p}\). Hence every
element outside \(\mathfrak pA_{\mathfrak p}\) is a unit. It follows that
\(\mathfrak pA_{\mathfrak p}\) contains all nonunits, so it is the unique
maximal ideal. Therefore \(A_{\mathfrak p}\) is local.

For \(A=k[x,y]\) and \(\mathfrak p=(x,y)\), a fraction
\[
\frac{f}{g}
\]
belongs to \(A_{\mathfrak p}\) exactly when
\[
g\notin(x,y),
\]
equivalently \(g(0,0)\ne0\). Such a fraction is a unit exactly when
\(f\notin(x,y)\), equivalently \(f(0,0)\ne0\).

For \(\mathbb Z_{(p)}\), the denominators are precisely the integers not
divisible by \(p\). A fraction \(a/b\) is a unit exactly when \(p\nmid a\).
Thus the nonunits are exactly the fractions whose numerator is divisible by
\(p\), namely
\[
p\mathbb Z_{(p)}.
\]
\end{solution}


\begin{problem}[CP-V-0017]
\label{prob:cp-v-0017}
For a commutative ring \(A\), prove that the following are equivalent.

\textbf{(i)} \(A\) is local.

\textbf{(ii)} The set of nonunits of \(A\) is an ideal.

\textbf{(iii)} There exists a proper ideal \(\mathfrak m\subset A\) such that
every element of \(A\setminus\mathfrak m\) is a unit.

When these conditions hold, prove that \(\mathfrak m\) is the unique maximal
ideal.

Use the criterion to verify directly that the following rings are local:
\[
k[[t]],\qquad k[x]_{(x)},\qquad \mathbb Z_{(p)}.
\]
\end{problem}

\begin{solution}
Assume \(A\) is local with unique maximal ideal \(\mathfrak m\). Every
nonunit lies in some maximal ideal, hence in \(\mathfrak m\). Conversely, no
element of \(\mathfrak m\) can be a unit. Thus the nonunits are exactly
\(\mathfrak m\), proving \textbf{(i)} implies \textbf{(ii)}.

If the nonunits form an ideal \(N\), then \(N\) is proper because \(1\) is a
unit. Every element outside \(N\) is a unit, so \textbf{(ii)} implies
\textbf{(iii)} with \(\mathfrak m=N\).

Assume \textbf{(iii)}. Any maximal ideal consists entirely of nonunits, hence
is contained in \(\mathfrak m\). Since \(\mathfrak m\) is proper, it is
contained in some maximal ideal. Therefore \(\mathfrak m\) itself is maximal
and every maximal ideal equals \(\mathfrak m\). Thus \(A\) is local.

For \(k[[t]]\), a power series is invertible exactly when its constant term is
nonzero. Hence the nonunits are exactly
\[
(t).
\]

For \(k[x]_{(x)}\), an element \(f/g\) with \(g(0)\ne0\) is a unit exactly
when \(f(0)\ne0\). Thus the nonunits form
\[
(x)k[x]_{(x)}.
\]

For \(\mathbb Z_{(p)}\), a fraction \(a/b\) with \(p\nmid b\) is a unit
exactly when \(p\nmid a\). Hence the nonunits form
\[
p\mathbb Z_{(p)}.
\]
Each ring is therefore local by the criterion.
\end{solution}


\begin{problem}[CP-V-0052]
\label{prob:cp-v-0052}
Let \(A\) be a ring and \(\mathfrak p\subseteq A\) a prime ideal.

\textbf{(a)} Prove that
\[
A_{\mathfrak p}/\mathfrak pA_{\mathfrak p}
\cong
\operatorname{Frac}(A/\mathfrak p).
\]

\textbf{(b)} If \(\mathfrak m\) is maximal, deduce that
\[
A/\mathfrak m
\cong
A_{\mathfrak m}/\mathfrak mA_{\mathfrak m}.
\]

\textbf{(c)} Prove that for every \(k\ge0\),
\[
\mathfrak m^k/\mathfrak m^{k+1}
\cong
(\mathfrak mA_{\mathfrak m})^k/
(\mathfrak mA_{\mathfrak m})^{k+1}.
\]
\end{problem}

\begin{solution}
Localizing the quotient gives
\[
A_{\mathfrak p}/\mathfrak pA_{\mathfrak p}
\cong
(A/\mathfrak p)_{A\setminus\mathfrak p}.
\]
The image of \(A\setminus\mathfrak p\) in the domain \(A/\mathfrak p\) is
exactly its set of nonzero elements, so the localization is the fraction
field.

If \(\mathfrak m\) is maximal, \(A/\mathfrak m\) is already a field, giving
the second isomorphism.

Localizing
\[
0\to\mathfrak m^{k+1}\to\mathfrak m^k
\to\mathfrak m^k/\mathfrak m^{k+1}\to0
\]
gives
\[
(\mathfrak m^k/\mathfrak m^{k+1})_{\mathfrak m}
\cong
(\mathfrak mA_{\mathfrak m})^k/
(\mathfrak mA_{\mathfrak m})^{k+1}.
\]
But \(\mathfrak m\) annihilates
\(\mathfrak m^k/\mathfrak m^{k+1}\), so every element outside
\(\mathfrak m\) acts invertibly on this \(A/\mathfrak m\)-vector space.
Thus localization does not change it.
\end{solution}


\begin{problem}[CP-V-0066]
\label{prob:cp-v-0066}
Determine whether
\[
\frac{1}{x}
\]
belongs to
\[
k[x,y]_{(x,y)}
\qquad\text{and}\qquad
k[x,y]_{(y)}.
\]
\end{problem}

\begin{solution}
In
\[
k[x,y]_{(x,y)},
\]
the denominator must not vanish at \((0,0)\). But
\[
x(0,0)=0,
\]
so
\[
x\in (x,y),
\]
and therefore
\[
\frac{1}{x}\notin k[x,y]_{(x,y)}.
\]

In
\[
k[x,y]_{(y)},
\]
the denominator is allowed if it is not divisible by \(y\). Since
\[
x\notin (y),
\]
we get
\[
\frac{1}{x}\in k[x,y]_{(y)}.
\]
\end{solution}


\begin{problem}[CP-V-0067]
\label{prob:cp-v-0067}
Determine whether
\[
\frac{1}{x+y}
\]
belongs to
\[
k[x,y]_{(x,y)}
\qquad\text{and}\qquad
k[x,y]_{(y)}.
\]
\end{problem}

\begin{solution}
At the origin,
\[
(x+y)(0,0)=0,
\]
so
\[
x+y\in (x,y).
\]
Hence
\[
\frac{1}{x+y}\notin k[x,y]_{(x,y)}.
\]

Now check whether
\[
x+y\in (y).
\]
This is false, because \(x+y\) is not divisible by \(y\). Therefore
\[
\frac{1}{x+y}\in k[x,y]_{(y)}.
\]
\end{solution}


\begin{problem}[CP-V-0068]
\label{prob:cp-v-0068}
Determine whether
\[
\frac{x^2+y}{1+x+y}
\]
belongs to
\[
k[x,y]_{(x,y)}.
\]
\end{problem}

\begin{solution}
We evaluate the denominator at the origin:
\[
(1+x+y)(0,0)=1.
\]
So
\[
1+x+y\notin (x,y).
\]
Hence
\[
\frac{x^2+y}{1+x+y}\in k[x,y]_{(x,y)}.
\]
\end{solution}


\begin{problem}[CP-V-0071]
\label{prob:cp-v-0071}
Determine whether
\[
\frac{1}{1+xy}
\]
belongs to
\[
k[x,y]_{(x,y)}.
\]
\end{problem}

\begin{solution}
We compute
\[
(1+xy)(0,0)=1.
\]
Thus
\[
1+xy\notin (x,y),
\]
so
\[
\frac{1}{1+xy}\in k[x,y]_{(x,y)}.
\]
\end{solution}


\begin{problem}[CP-V-0072]
\label{prob:cp-v-0072}
Determine whether
\[
\frac{1}{y^2+x}
\]
belongs to
\[
k[x,y]_{(y)}.
\]
\end{problem}

\begin{solution}
To belong to
\[
k[x,y]_{(y)},
\]
the denominator must not lie in \((y)\). The polynomial
\[
y^2+x
\]
is not divisible by \(y\), since it has the term \(x\). Therefore
\[
y^2+x\notin (y),
\]
and hence
\[
\frac{1}{y^2+x}\in k[x,y]_{(y)}.
\]
\end{solution}


\begin{problem}[CP-V-0075]
\label{prob:cp-v-0075}
Let
\[
A=k[x],
\qquad
\mathfrak p=(x-a).
\]
Show that
\[
\frac{f(x)}{g(x)}\in A_{\mathfrak p}
\]
if and only if
\[
g(a)\neq 0.
\]
\end{problem}

\begin{solution}
By definition,
\[
A_{\mathfrak p}=S^{-1}A,
\qquad
S=A\setminus \mathfrak p.
\]
So \(\frac{f(x)}{g(x)}\) belongs to \(A_{\mathfrak p}\) exactly when
\[
g(x)\notin \mathfrak p=(x-a).
\]
By the factor theorem,
\[
g(x)\in (x-a)
\quad\Longleftrightarrow\quad
g(a)=0.
\]
Therefore
\[
g(x)\notin (x-a)
\quad\Longleftrightarrow\quad
g(a)\neq 0.
\]
Hence
\[
\frac{f(x)}{g(x)}\in k[x]_{(x-a)}
\quad\Longleftrightarrow\quad
g(a)\neq 0.
\]
\end{solution}


\begin{problem}[CP-V-0076]
\label{prob:cp-v-0076}
Let
\[
A=k[x,y],
\qquad
\mathfrak m=(x-a,y-b).
\]
Show that
\[
\frac{f(x,y)}{g(x,y)}\in A_{\mathfrak m}
\]
if and only if
\[
g(a,b)\neq 0.
\]
\end{problem}

\begin{solution}
By definition,
\[
A_{\mathfrak m}=S^{-1}A,
\qquad
S=A\setminus \mathfrak m.
\]
Thus \(\frac{f}{g}\) belongs to \(A_{\mathfrak m}\) exactly when
\[
g\notin \mathfrak m.
\]
But
\[
\mathfrak m=(x-a,y-b)
\]
consists exactly of those polynomials vanishing at \((a,b)\). Therefore
\[
g\notin (x-a,y-b)
\quad\Longleftrightarrow\quad
g(a,b)\neq 0.
\]
So
\[
\frac{f(x,y)}{g(x,y)}\in k[x,y]_{(x-a,y-b)}
\quad\Longleftrightarrow\quad
g(a,b)\neq 0.
\]
\end{solution}


\begin{problem}[CP-V-0099]
\label{prob:cp-v-0099}
Compute \(k[t]_{(t-a)}\) and describe its units.
\end{problem}

\begin{solution}
The localization is
\[
k[t]_{(t-a)}
=
\left\{\frac{f(t)}{g(t)}:g(a)\ne0\right\}.
\]
Its maximal ideal is generated by \(t-a\). A fraction \(f/g\) is a unit exactly when \(f(a)\ne0\).
\end{solution}


\begin{problem}[CP-V-0100]
\label{prob:cp-v-0100}
Compute the local ring of \(k[s,t]/(st)\) at \((s,t)\), and compare it with local rings at points on one component away from the intersection.
\end{problem}

\begin{solution}
At the crossing,
\[
A_{(s,t)}=\bigl(k[s,t]/(st)\bigr)_{(s,t)}
\]
has two minimal primes \((s)\) and \((t)\), so both branches remain visible. At a point \((s-a,t)\) with \(a\ne0\), the element \(s\) is a unit; from \(st=0\) it follows that \(t=0\), so the local ring becomes the regular local ring of the \(s\)-axis at that point. The analogous statement holds on the other component.
\end{solution}



\section*{V/09 --- Modules and Exact Sequences}

\begin{problem}[CP-V-0024]
\label{prob:cp-v-0024}
Let \(R\) be a ring.

\textbf{(a)} Prove that the functor
\[
-\otimes_R N
\]
is right exact for every \(R\)-module \(N\).

\textbf{(b)} Show that tensor product need not be left exact by tensoring
\[
0\longrightarrow\mathbb Z
\xrightarrow{\cdot2}
\mathbb Z
\longrightarrow
\mathbb Z/2\mathbb Z
\longrightarrow0
\]
with \(\mathbb Z/2\mathbb Z\).

\textbf{(c)} Prove that
\[
\operatorname{Hom}_R(M,-)
\]
is left exact for every \(R\)-module \(M\).

\textbf{(d)} Show that Hom need not be right exact by applying
\[
\operatorname{Hom}_{\mathbb Z}(\mathbb Z/2\mathbb Z,-)
\]
to the same short exact sequence.
\end{problem}

\begin{solution}
Tensor product is a left adjoint, hence preserves cokernels. Concretely, if
\[
A'\longrightarrow A\longrightarrow A''\longrightarrow0
\]
is exact, then
\[
A'\otimes_RN
\longrightarrow
A\otimes_RN
\longrightarrow
A''\otimes_RN
\longrightarrow0
\]
is exact.

For the displayed sequence over \(\mathbb Z\), tensoring with
\(\mathbb Z/2\mathbb Z\) gives
\[
0\longrightarrow
\mathbb Z/2\mathbb Z
\xrightarrow{0}
\mathbb Z/2\mathbb Z
\longrightarrow
\mathbb Z/2\mathbb Z
\longrightarrow0.
\]
The first map is not injective, so tensor product is not left exact in
general.

For Hom, if
\[
0\longrightarrow A'\longrightarrow A\longrightarrow A''
\]
is exact, then
\[
0\longrightarrow
\operatorname{Hom}_R(M,A')
\longrightarrow
\operatorname{Hom}_R(M,A)
\longrightarrow
\operatorname{Hom}_R(M,A'')
\]
is exact. Indeed, a map \(M\to A\) lands in \(A'\) exactly when its composite
with \(A\to A''\) is zero.

Finally,
\[
\operatorname{Hom}_{\mathbb Z}(\mathbb Z/2\mathbb Z,\mathbb Z)=0,
\]
whereas
\[
\operatorname{Hom}_{\mathbb Z}
(\mathbb Z/2\mathbb Z,\mathbb Z/2\mathbb Z)
\cong
\mathbb Z/2\mathbb Z.
\]
Applying Hom to the displayed short exact sequence therefore ends with
\[
0\longrightarrow0\longrightarrow0
\longrightarrow
\mathbb Z/2\mathbb Z,
\]
so the final map is not surjective. Thus Hom is not right exact in general.
\end{solution}


\begin{problem}[CP-V-0049]
\label{prob:cp-v-0049}
Let \(A\) be an integral domain and \(M\) an \(A\)-module. Let
\[
T(M)=\{m\in M:\exists\,0\ne a\in A,\ am=0\}.
\]

Prove that:

\textbf{(a)} \(T(M)\) is a submodule of \(M\);

\textbf{(b)} \(M/T(M)\) is torsion-free;

\textbf{(c)} every homomorphism \(f:M\to N\) satisfies
\[
f(T(M))\subseteq T(N);
\]

\textbf{(d)} for every multiplicatively closed \(S\subseteq A\),
\[
T(S^{-1}M)=S^{-1}T(M).
\]
\end{problem}

\begin{solution}
If \(m,n\in T(M)\), choose nonzero \(a,b\in A\) with \(am=0\) and \(bn=0\).
Since \(A\) is a domain, \(ab\ne0\), and
\[
ab(m+n)=0.
\]
Also \(a(rm)=r(am)=0\), so \(T(M)\) is a submodule.

If \(a(m+T(M))=0\) in \(M/T(M)\), then \(am\in T(M)\), so some nonzero
\(b\) satisfies \(bam=0\). Since \(ba\ne0\), \(m\in T(M)\). Thus the quotient
is torsion-free.

If \(am=0\), then
\[
a f(m)=f(am)=0,
\]
so homomorphisms preserve torsion.

Finally, if \(m/s\in T(S^{-1}M)\), then for some nonzero \(a/t\in S^{-1}A\),
\[
\frac{a}{t}\frac{m}{s}=0.
\]
Hence for some \(u\in S\),
\[
uam=0.
\]
Because \(ua\ne0\), the element \(m\) is torsion, so
\[
m/s\in S^{-1}T(M).
\]
The reverse inclusion is immediate.
\end{solution}



\section*{V/10 --- Tensor Products}

\begin{problem}[CP-V-0018]
\label{prob:cp-v-0018}
Let \(k\) be a field.

\textbf{(a)} Compute
\[
k\otimes_{k[x,y]}k[u,v],
\]
where \(k[x,y]\to k\) sends \(x,y\mapsto0\), while
\(k[x,y]\to k[u,v]\) sends
\[
x\mapsto u,\qquad y\mapsto uv.
\]

\textbf{(b)} Compute
\[
k[v]\otimes_{k[x]}k[v],
\]
where both copies are \(k[x]\)-algebras via \(x\mapsto v^2\).
Describe separately the cases \(\operatorname{char}k\ne2\) and
\(\operatorname{char}k=2\).
\end{problem}

\begin{solution}
For \textbf{(a)}, since
\[
k\cong k[x,y]/(x,y),
\]
base change gives
\[
k\otimes_{k[x,y]}k[u,v]
\cong
k[u,v]/(x,y)k[u,v].
\]
Under the given algebra map,
\[
(x,y)k[u,v]=(u,uv)=(u),
\]
so
\[
k\otimes_{k[x,y]}k[u,v]
\cong
k[u,v]/(u)
\cong
k[v].
\]

For \textbf{(b)}, distinguish the two polynomial variables as \(v\) and \(w\).
Then
\[
k[v]\otimes_{k[x]}k[w]
\cong
k[v,w]/(v^2-w^2).
\]
If \(\operatorname{char}k\ne2\), then
\[
v^2-w^2=(v-w)(v+w),
\]
and the two factors generate comaximal ideals. Hence the Chinese remainder
theorem gives
\[
k[v,w]/(v^2-w^2)
\cong
k[v,w]/(v-w)\times k[v,w]/(v+w)
\cong
k[v]\times k[v].
\]

If \(\operatorname{char}k=2\), then
\[
v^2-w^2=(v-w)^2,
\]
so
\[
k[v]\otimes_{k[x]}k[w]
\cong
k[v,w]/(v-w)^2,
\]
which is nonreduced.
\end{solution}


\begin{problem}[CP-V-0095]
\label{prob:cp-v-0095}
Compute \(\mathbb C\otimes_{\mathbb R}\mathbb C\).
\end{problem}

\begin{solution}
Using \(\mathbb C\cong\mathbb R[t]/(t^2+1)\),
\[
\mathbb C\otimes_{\mathbb R}\mathbb C
\cong
\mathbb C[t]/(t^2+1)
=
\mathbb C[t]/((t-i)(t+i)).
\]
The two factors are comaximal, so the Chinese remainder theorem gives
\[
\mathbb C\otimes_{\mathbb R}\mathbb C\cong\mathbb C\times\mathbb C.
\]
\end{solution}



\section*{V/11 --- Quotients and Base Change}

\begin{problem}[CP-V-0022]
\label{prob:cp-v-0022}
Let \(A\) be a ring, let \(I\subseteq A\) be an ideal, and let \(M\) be an
\(A\)-module. Prove that there is a natural isomorphism
\[
(A/I)\otimes_A M \cong M/IM.
\]
Describe explicitly the mutually inverse maps.
\end{problem}

\begin{solution}
Define
\[
\Psi:(A/I)\otimes_A M\longrightarrow M/IM
\]
on pure tensors by
\[
\Psi((a+I)\otimes m)=am+IM.
\]
If \(a-a'\in I\), then
\[
(a-a')m\in IM,
\]
so the definition does not depend on the chosen representative of \(a+I\).
The map is \(A\)-balanced and therefore induces a homomorphism on the tensor
product.

It is surjective because
\[
\Psi((1+I)\otimes m)=m+IM.
\]

Conversely, define
\[
\Theta:M\longrightarrow (A/I)\otimes_A M,
\qquad
\Theta(m)=(1+I)\otimes m.
\]
If \(m\in IM\), write
\[
m=\sum_j i_jm_j,\qquad i_j\in I.
\]
Then
\[
\Theta(m)
=
\sum_j (1+I)\otimes i_jm_j
=
\sum_j(i_j+I)\otimes m_j
=
0.
\]
Hence \(\Theta\) factors through \(M/IM\):
\[
\overline\Theta:M/IM\longrightarrow(A/I)\otimes_A M.
\]
Now
\[
\Psi\circ\overline\Theta=\operatorname{id}_{M/IM},
\]
and on pure tensors,
\[
\overline\Theta\bigl(\Psi((a+I)\otimes m)\bigr)
=
(1+I)\otimes am
=
(a+I)\otimes m.
\]
Thus the maps are inverse, proving
\[
(A/I)\otimes_A M\cong M/IM.
\]
\end{solution}


\begin{problem}[CP-V-0096]
\label{prob:cp-v-0096}
Show that \((k[x]/(f))\otimes_k K\cong K[x]/(f)\) for a field extension \(K/k\).
\end{problem}

\begin{solution}
Tensor the exact presentation
\[
k[x]\xrightarrow{\ \pi\ }k[x]/(f)\to0
\]
with \(K\). Since \(K\) is flat over \(k\),
\[
K\otimes_k(k[x]/(f))
\cong
(K\otimes_k k[x])/(f)
\cong
K[x]/(f).
\]
\end{solution}



\section*{V/12 --- Hom and Finitely Presented Modules}

\begin{problem}[CP-V-0023]
\label{prob:cp-v-0023}
Let \(A\) be a ring, \(S\subseteq A\) a multiplicatively closed set, \(N\) an
\(A\)-module, and \(M\) a finitely presented \(A\)-module. Prove that
\[
S^{-1}\!\operatorname{Hom}_A(M,N)
\cong
\operatorname{Hom}_{S^{-1}A}(S^{-1}M,S^{-1}N).
\]
Give the explicit formula for the canonical map.
\end{problem}

\begin{solution}
Choose a finite presentation
\[
A^{n_1}\xrightarrow{u}A^{n_2}\longrightarrow M\longrightarrow0.
\]
Applying the contravariant functor \(\operatorname{Hom}_A(-,N)\) gives
\[
0\longrightarrow\operatorname{Hom}_A(M,N)
\longrightarrow N^{n_2}
\xrightarrow{u^*}
N^{n_1}.
\]
Localization is exact and commutes with finite direct sums, so
\[
0\longrightarrow
S^{-1}\operatorname{Hom}_A(M,N)
\longrightarrow
(S^{-1}N)^{n_2}
\longrightarrow
(S^{-1}N)^{n_1}.
\]

On the other hand, localizing the presentation of \(M\) gives
\[
(S^{-1}A)^{n_1}
\longrightarrow
(S^{-1}A)^{n_2}
\longrightarrow
S^{-1}M
\longrightarrow0.
\]
Applying
\[
\operatorname{Hom}_{S^{-1}A}(-,S^{-1}N)
\]
produces
\[
0\longrightarrow
\operatorname{Hom}_{S^{-1}A}(S^{-1}M,S^{-1}N)
\longrightarrow
(S^{-1}N)^{n_2}
\longrightarrow
(S^{-1}N)^{n_1}.
\]
The two left-hand modules are therefore kernels of the same map, so they are
canonically isomorphic.

Explicitly,
\[
\frac{f}{s}
\longmapsto
\left(
\frac{m}{t}\longmapsto\frac{f(m)}{st}
\right).
\]
The finite-presentation hypothesis is what allows the comparison with finite
free modules and finite direct sums.
\end{solution}


\begin{problem}[CP-V-0105]
\label{prob:cp-v-0105}
Let \(M\) be a finitely presented \(A\)-module and let \((N_i)_{i\in I}\) be a filtered direct system of \(A\)-modules. Prove that the canonical map
\[
\varinjlim_i \operatorname{Hom}_A(M,N_i)
\longrightarrow
\operatorname{Hom}_A\!\left(M,\varinjlim_i N_i\right)
\]
is an isomorphism.
\end{problem}

\begin{solution}
Choose a finite presentation
\[
A^r\longrightarrow A^s\longrightarrow M\longrightarrow0.
\]
Applying \(\operatorname{Hom}_A(-,N_i)\) gives
\[
0\longrightarrow\operatorname{Hom}_A(M,N_i)
\longrightarrow N_i^s\longrightarrow N_i^r.
\]
Filtered colimits are exact and commute with finite direct sums, so after taking \(\varinjlim_i\) we obtain
\[
0\longrightarrow\varinjlim_i\operatorname{Hom}_A(M,N_i)
\longrightarrow \left(\varinjlim_iN_i\right)^s
\longrightarrow \left(\varinjlim_iN_i\right)^r.
\]
Applying \(\operatorname{Hom}_A(M,-)\) directly to the same presentation with \(N=\varinjlim_iN_i\) identifies
\[
\operatorname{Hom}_A\!\left(M,\varinjlim_iN_i\right)
\]
with the same kernel. Hence the canonical map is an isomorphism.
\end{solution}



\section*{V/13 --- Free and Projective Modules}

\begin{problem}[CP-V-0020]
\label{prob:cp-v-0020}
Let
\[
A=A_1\times A_2,
\qquad
M_1=A_1\times0,
\qquad
M_2=0\times A_2,
\]
where \(A_1,A_2\ne0\).

\textbf{(a)} Prove that \(M_1\) and \(M_2\) are projective \(A\)-modules.

\textbf{(b)} Show that neither \(M_1\) nor \(M_2\) is free.

\textbf{(c)} Interpret the decomposition using the idempotents
\[
e_1=(1,0),\qquad e_2=(0,1).
\]
\end{problem}

\begin{solution}
We have a direct-sum decomposition
\[
A=(A_1\times0)\oplus(0\times A_2)=M_1\oplus M_2.
\]
Since \(A\) is a free \(A\)-module of rank one, each \(M_i\) is a direct
summand of a free module. Therefore \(M_1\) and \(M_2\) are projective.

To show that \(M_1\) is not free, compute its annihilator:
\[
\operatorname{Ann}_A(M_1)=0\times A_2.
\]
This is nonzero because \(A_2\ne0\). But every nonzero free \(A\)-module has
zero annihilator: if \(aA^n=0\), then applying \(a\) to a standard basis
vector gives \(a=0\). Hence \(M_1\) is not free. Similarly,
\[
\operatorname{Ann}_A(M_2)=A_1\times0\ne0,
\]
so \(M_2\) is not free.

Finally,
\[
e_1^2=e_1,\qquad e_2^2=e_2,\qquad e_1e_2=0,\qquad e_1+e_2=1,
\]
and
\[
M_1=Ae_1,\qquad M_2=Ae_2.
\]
Thus ideals generated by complementary idempotents give projective direct
summands of \(A\).
\end{solution}


\begin{problem}[CP-V-0021]
\label{prob:cp-v-0021}
Let \((A,\mathfrak m)\) be a local ring, let
\[
k=A/\mathfrak m,
\]
and let \(M\) be a finitely generated \(A\)-module.

\textbf{(a)} Let \(f:F\to M\) be a homomorphism with \(F\) finitely generated
and free. Prove that the following are equivalent:
\[
f\otimes_A k:F\otimes_A k\longrightarrow M\otimes_A k
\]
is an isomorphism, and
\[
f\text{ is surjective with }\ker f\subseteq\mathfrak m F.
\]

\textbf{(b)} Prove that there exists such a map \(f:F\to M\).

\textbf{(c)} Show that the rank of \(F\) in any such minimal free cover is
\[
\dim_k(M/\mathfrak mM).
\]
\end{problem}

\begin{solution}
Using
\[
L\otimes_A k\cong L/\mathfrak mL
\]
for any \(A\)-module \(L\), the map \(f\otimes_A k\) identifies with
\[
\bar f:F/\mathfrak mF\longrightarrow M/\mathfrak mM.
\]

Assume first that \(\bar f\) is an isomorphism. Surjectivity of \(\bar f\)
gives
\[
M=f(F)+\mathfrak mM.
\]
Since \(M\) is finitely generated, Nakayama's lemma yields
\[
M=f(F),
\]
so \(f\) is surjective.

If \(x\in\ker f\), then the class of \(x\) maps to zero under \(\bar f\).
Because \(\bar f\) is injective,
\[
x\in\mathfrak mF.
\]
Hence
\[
\ker f\subseteq\mathfrak mF.
\]

Conversely, suppose \(f\) is surjective and
\[
\ker f\subseteq\mathfrak mF.
\]
Then \(\bar f\) is surjective. If the class of \(x\in F\) belongs to
\(\ker\bar f\), then
\[
f(x)\in\mathfrak mM.
\]
Since \(f\) is surjective,
\[
\mathfrak mM=f(\mathfrak mF),
\]
so there is \(y\in\mathfrak mF\) with \(f(x)=f(y)\). Thus
\[
x-y\in\ker f\subseteq\mathfrak mF,
\]
and therefore \(x\in\mathfrak mF\). Hence \(\bar f\) is injective.

For existence, choose elements
\[
m_1,\dots,m_r\in M
\]
whose residue classes form a \(k\)-basis of
\[
M/\mathfrak mM.
\]
Let \(F=A^r\) and define
\[
f:F\to M
\]
by sending the \(i\)-th basis vector to \(m_i\). Modulo \(\mathfrak m\), this
map is an isomorphism by construction, so part \textbf{(a)} applies.

Finally,
\[
F/\mathfrak mF\cong k^{\operatorname{rank}F}.
\]
Since a minimal free cover induces an isomorphism
\[
F/\mathfrak mF\cong M/\mathfrak mM,
\]
we obtain
\[
\operatorname{rank}F
=
\dim_k(M/\mathfrak mM).
\]
\end{solution}


\begin{problem}[CP-V-0025]
\label{prob:cp-v-0025}
Let \(R\) be a ring and \(P\) an \(R\)-module. Prove that the following are
equivalent.

\textbf{(i)} \(P\) is projective.

\textbf{(ii)} Every surjection
\[
E\twoheadrightarrow P
\]
splits.

\textbf{(iii)} Every short exact sequence
\[
0\longrightarrow K\longrightarrow E\longrightarrow P\longrightarrow0
\]
splits.

Use this criterion to prove that
\[
R/(x)
\]
is not projective over \(R=k[x]\).
\end{problem}

\begin{solution}
Assume \(P\) is projective and let
\[
\pi:E\twoheadrightarrow P
\]
be surjective. Apply the lifting property to
\[
\operatorname{id}_P:P\to P.
\]
There is a map
\[
s:P\to E
\]
with
\[
\pi\circ s=\operatorname{id}_P.
\]
Thus \(\pi\) splits, proving \textbf{(i)} implies \textbf{(ii)}.

Statement \textbf{(ii)} immediately implies \textbf{(iii)}, since every short
exact sequence ending in \(P\) contains such a surjection.

Conversely, suppose every short exact sequence ending in \(P\) splits. Choose
a surjection from a free module
\[
F\twoheadrightarrow P.
\]
Its kernel \(K\) gives
\[
0\longrightarrow K\longrightarrow F\longrightarrow P\longrightarrow0.
\]
By assumption this sequence splits, so
\[
F\cong K\oplus P.
\]
Hence \(P\) is a direct summand of a free module and is therefore projective.

Now let
\[
R=k[x].
\]
Consider
\[
0\longrightarrow R
\xrightarrow{\cdot x}
R
\longrightarrow
R/(x)
\longrightarrow0.
\]
If this sequence split, then \(R/(x)\) would be a direct summand of \(R\).
Equivalently, the ideal \((x)\) would be generated by an idempotent. But the
only idempotents in the domain \(k[x]\) are \(0\) and \(1\), and \((x)\) is
neither \(0\) nor \(R\). Hence the sequence does not split, and \(R/(x)\) is
not projective.
\end{solution}


\begin{problem}[CP-V-0058]
\label{prob:cp-v-0058}
Let \(A\) be a ring and let \(F=A^n\). Suppose
\[
x_1,\dots,x_n
\]
generate \(F\). Prove that they form a basis.
\end{problem}

\begin{solution}
Define
\[
\varphi:A^n\to A^n,\qquad e_i\mapsto x_i.
\]
Since the \(x_i\) generate, \(\varphi\) is surjective.

Localize at a maximal ideal \(\mathfrak m\). Then
\[
\varphi_{\mathfrak m}:A_{\mathfrak m}^n\to A_{\mathfrak m}^n
\]
is surjective. Reducing modulo the maximal ideal gives a surjective linear
map
\[
k(\mathfrak m)^n\to k(\mathfrak m)^n,
\]
hence an isomorphism. Therefore the determinant of a matrix representing
\(\varphi_{\mathfrak m}\) is a unit, so \(\varphi_{\mathfrak m}\) is an
isomorphism.

Thus
\[
(\ker\varphi)_{\mathfrak m}=0
\]
for every maximal ideal \(\mathfrak m\). A module whose localization at every
maximal ideal is zero is itself zero, so
\[
\ker\varphi=0.
\]
Hence \(\varphi\) is an isomorphism and \(x_1,\dots,x_n\) form a basis.
\end{solution}



\section*{V/14 --- Flat Modules}

\begin{problem}[CP-V-0019]
\label{prob:cp-v-0019}
Let \(A\) be a ring.

\textbf{(a)} Prove that \(A[x]\) is a free \(A\)-module.

\textbf{(b)} Deduce that \(A[x]\) is flat over \(A\).

\textbf{(c)} More generally, prove that
\[
A[x_1,\dots,x_n]
\]
is flat over \(A\).

\textbf{(d)} Explain why this gives an elementary example of a faithfully flat
extension \(A\to A[x]\).
\end{problem}

\begin{solution}
As an \(A\)-module,
\[
A[x]
=
\bigoplus_{n\ge0}A x^n.
\]
Thus \(A[x]\) is free with basis
\[
1,x,x^2,\dots.
\]
Every free module is flat, so \(A[x]\) is flat over \(A\).

Likewise,
\[
A[x_1,\dots,x_n]
\]
is free over \(A\), with basis consisting of all monomials
\[
x_1^{a_1}\cdots x_n^{a_n},
\qquad a_i\ge0.
\]
Hence it is flat.

Finally, the inclusion \(A\to A[x]\) admits the \(A\)-linear retraction
\[
A[x]\longrightarrow A,\qquad f(x)\longmapsto f(0).
\]
Thus \(A\) is a direct summand of \(A[x]\) as an \(A\)-module. If
\(M\otimes_A A[x]=0\), applying the induced retraction shows
\[
M\cong M\otimes_A A=0.
\]
Therefore tensoring with \(A[x]\) is faithful as well as exact, so
\(A\to A[x]\) is faithfully flat.
\end{solution}


\begin{problem}[CP-V-0106]
\label{prob:cp-v-0106}
Let \(A\) be a principal ideal domain and \(M\) an \(A\)-module. Prove that \(M\) is flat if and only if \(M\) is torsion-free.
\end{problem}

\begin{solution}
Every flat module over a domain is torsion-free: if \(0\ne a\in A\), the injection
\[
A\xrightarrow{\cdot a}A
\]
remains injective after tensoring with \(M\), so multiplication by \(a\) on \(M\) is injective.

Conversely, assume \(M\) is torsion-free. Over a PID, every finitely generated submodule \(N\subseteq M\) is free. Thus \(M\) is the filtered union of its finitely generated free submodules. Equivalently,
\[
M\cong\varinjlim N.
\]
Free modules are flat, and filtered colimits of flat modules are flat. Therefore \(M\) is flat.
\end{solution}



\section*{V/15 --- Noetherian Rings and Modules}

\begin{problem}[CP-V-0050]
\label{prob:cp-v-0050}
Let \(A\) be a ring such that:

\textbf{(i)} \(A_{\mathfrak m}\) is Noetherian for every maximal ideal
\(\mathfrak m\subseteq A\);

\textbf{(ii)} for every nonzero \(x\in A\), only finitely many maximal ideals
contain \(x\).

Prove that \(A\) is Noetherian.
\end{problem}

\begin{solution}
Let \(I\ne0\) be an ideal and choose \(0\ne x\in I\). By hypothesis, \(x\)
lies in only finitely many maximal ideals
\[
\mathfrak m_1,\dots,\mathfrak m_r.
\]
Since each \(A_{\mathfrak m_i}\) is Noetherian, choose finitely many elements
of \(I\) whose images generate \(I_{\mathfrak m_i}\). Let \(J\subseteq I\)
be the ideal generated by \(x\) and all these finitely many chosen elements.

For a maximal ideal \(\mathfrak m\), if \(x\notin\mathfrak m\), then
\(x/1\) is a unit in \(A_{\mathfrak m}\), so
\[
J_{\mathfrak m}=A_{\mathfrak m}=I_{\mathfrak m}.
\]
If \(x\in\mathfrak m\), then \(\mathfrak m=\mathfrak m_i\) for some \(i\),
and by construction
\[
J_{\mathfrak m}=I_{\mathfrak m}.
\]
Thus
\[
(I/J)_{\mathfrak m}=0
\]
for every maximal ideal \(\mathfrak m\). Hence \(I/J=0\), so \(I=J\) is
finitely generated. Therefore every ideal of \(A\) is finitely generated,
and \(A\) is Noetherian.
\end{solution}


\begin{problem}[CP-V-0056]
\label{prob:cp-v-0056}
Let
\[
S=k[x^3,x^2y,xy^2,y^3]\subseteq k[x,y]
\]
and
\[
\mathfrak m=(x^3,x^2y,xy^2,y^3).
\]
Prove that the local ring \(S_{\mathfrak m}\) is Cohen--Macaulay by showing
that
\[
x^3,\ y^3
\]
is a regular sequence.
\end{problem}

\begin{solution}
Let
\[
R=k[x^3,y^3]\cong k[u,v].
\]
Every monomial of \(S\) has total degree divisible by \(3\), and according
to the exponent of \(x\) modulo \(3\),
\[
S=R\oplus R\,x^2y\oplus R\,xy^2.
\]
After localizing at
\[
\mathfrak n=(x^3,y^3)\subseteq R,
\]
we obtain
\[
S_{\mathfrak m}
=
R_{\mathfrak n}
\oplus
R_{\mathfrak n}x^2y
\oplus
R_{\mathfrak n}xy^2.
\]
Thus \(S_{\mathfrak m}\) is free over the \(2\)-dimensional regular local
domain \(R_{\mathfrak n}\).

Multiplication by \(x^3\) is injective on \(R_{\mathfrak n}\), hence on the
free module \(S_{\mathfrak m}\). Modulo \(x^3\),
\[
R_{\mathfrak n}/(x^3)
\cong
k[y^3]_{(y^3)}
\]
is a domain, so multiplication by \(y^3\) remains injective on each of the
three summands. Therefore
\[
x^3,y^3
\]
is an \(S_{\mathfrak m}\)-regular sequence.

Since
\[
\dim S_{\mathfrak m}=2,
\]
we have
\[
\operatorname{depth}S_{\mathfrak m}=2=\dim S_{\mathfrak m},
\]
so \(S_{\mathfrak m}\) is Cohen--Macaulay.
\end{solution}


\begin{problem}[CP-V-0057]
\label{prob:cp-v-0057}
Let
\[
A=k[x_1,x_2,x_3,\dots],
\]
and choose integers
\[
m_1<m_2<\cdots
\]
whose gaps
\[
d_i=m_{i+1}-m_i
\]
are strictly increasing. Put
\[
\mathfrak p_i=(x_{m_i+1},\dots,x_{m_{i+1}})
\]
and
\[
S=A\setminus\bigcup_{i\ge1}\mathfrak p_i.
\]

Prove that:

\textbf{(a)} \(S\) is multiplicatively closed;

\textbf{(b)} the maximal ideals of \(S^{-1}A\) are exactly
\[
S^{-1}\mathfrak p_i;
\]

\textbf{(c)} \(S^{-1}A\) is Noetherian;

\textbf{(d)}
\[
\operatorname{ht}(S^{-1}\mathfrak p_i)=d_i,
\]
and therefore
\[
\dim S^{-1}A=\infty.
\]
\end{problem}

\begin{solution}
Each \(\mathfrak p_i\) is generated by a block of variables disjoint from the
blocks defining the other \(\mathfrak p_j\). Suppose \(fg\notin S\). Then
\(fg\in\mathfrak p_i\) for some \(i\). Since \(\mathfrak p_i\) is prime,
\(f\in\mathfrak p_i\) or \(g\in\mathfrak p_i\), so \(f\notin S\) or
\(g\notin S\). Hence \(S\) is multiplicatively closed.

A prime of \(S^{-1}A\) corresponds to a prime \(\mathfrak q\subseteq A\)
disjoint from \(S\), equivalently
\[
\mathfrak q\subseteq\bigcup_i\mathfrak p_i.
\]
By prime avoidance for this disjoint block family, such a prime is contained
in some \(\mathfrak p_i\). Therefore the maximal primes disjoint from \(S\)
are precisely the \(\mathfrak p_i\), and the maximal ideals of the localized
ring are
\[
S^{-1}\mathfrak p_i.
\]

At \(S^{-1}\mathfrak p_i\), all variables outside the finite block
\[
x_{m_i+1},\dots,x_{m_{i+1}}
\]
become units. Thus the local ring is obtained from a polynomial ring in the
\(d_i\) block variables over a field of fractions in the remaining
variables. In particular, each localization at a maximal ideal is
Noetherian. Moreover, every nonzero element of \(S^{-1}A\) lies in only
finitely many maximal ideals, because a polynomial involves only finitely
many variables. By the local-global criterion of CP-V-0050,
\[
S^{-1}A
\]
is Noetherian.

The localized maximal ideal is generated by the \(d_i\) block variables and
has a chain of prime ideals of length \(d_i\). Hence
\[
\operatorname{ht}(S^{-1}\mathfrak p_i)=d_i.
\]
Since the \(d_i\) are unbounded,
\[
\dim S^{-1}A=\infty.
\]
\end{solution}


\begin{problem}[CP-V-0083]
\label{prob:cp-v-0083}
Let \(A\) be a ring of finite Krull dimension. Prove
\[
1+\dim A\le \dim A[x]\le 1+2\dim A.
\]
\end{problem}

\begin{solution}
For the lower bound, a chain
\[
\mathfrak p_0\subsetneq\cdots\subsetneq\mathfrak p_d
\]
in \(A\) gives
\[
\mathfrak p_0A[x]\subsetneq\cdots\subsetneq
\mathfrak p_dA[x]\subsetneq(\mathfrak p_d,x),
\]
so
\[
\dim A[x]\ge d+1.
\]

For the upper bound, take a chain of primes in \(A[x]\) and contract it to
\(A\). For any fixed contracted prime \(\mathfrak p\), the primes lying over
\(\mathfrak p\) correspond after quotienting and localizing to primes of
\[
\operatorname{Frac}(A/\mathfrak p)[x],
\]
whose dimension is \(1\). Thus between two strict increases of contracted
primes there can be at most one extra strict inclusion in the fiber. A chain
whose contractions have length at most \(\dim A\) therefore has length at
most
\[
2\dim A+1.
\]
Hence
\[
\dim A[x]\le 1+2\dim A.
\]
\end{solution}


\begin{problem}[CP-V-0084]
\label{prob:cp-v-0084}
Let
\[
S=k[x_0,\dots,x_n]
\]
with its standard grading.

\textbf{(a)} If \(f\in S_d\) is homogeneous and nonzero, compute the Hilbert
polynomial of \(S/(f)\).

\textbf{(b)} If \(f\in S_d\) and \(g\in S_e\) form a regular sequence,
compute the Hilbert polynomial of \(S/(f,g)\).
\end{problem}

\begin{solution}
Since
\[
0\to S(-d)\xrightarrow{\cdot f}S\to S/(f)\to0,
\]
we obtain
\[
P_{S/(f)}(t)
=
\binom{t+n}{n}
-
\binom{t-d+n}{n}.
\]

If \(g\) is a non-zero-divisor on \(S/(f)\), then
\[
0\to S/(f)(-e)\xrightarrow{\cdot g}S/(f)
\to S/(f,g)\to0.
\]
Therefore
\[
P_{S/(f,g)}(t)
=
P_{S/(f)}(t)-P_{S/(f)}(t-e),
\]
hence
\[
P_{S/(f,g)}(t)
=
\binom{t+n}{n}
-\binom{t-d+n}{n}
-\binom{t-e+n}{n}
+\binom{t-d-e+n}{n}.
\]
The leading terms show that the degrees are \(n-1\) and \(n-2\),
respectively.
\end{solution}


\begin{problem}[CP-V-0085]
\label{prob:cp-v-0085}
Let \(I\subseteq S=k[x_0,\dots,x_n]\) be homogeneous, and define the degree
of \(S/I\) from the leading coefficient of its Hilbert polynomial.

\textbf{(a)} Show that for a hypersurface of degree \(d\),
\[
\deg S/(f)=d.
\]

\textbf{(b)} If \(f,g\) form a homogeneous regular sequence of degrees
\(d,e\), show
\[
\deg S/(f,g)=de.
\]

\textbf{(c)} More generally, if \(h\) is a homogeneous non-zero-divisor on
\(S/I\) of degree \(e\), prove
\[
\deg S/(I,h)=e\,\deg S/I.
\]
\end{problem}

\begin{solution}
The first two assertions follow from the leading terms of the Hilbert
polynomials computed in CP-V-0084.

For \textbf{(c)}, if
\[
P_{S/I}(t)=a_rt^r+\text{lower terms},
\]
then the exact sequence
\[
0\to S/I(-e)\xrightarrow{\cdot h}S/I
\to S/(I,h)\to0
\]
gives
\[
P_{S/(I,h)}(t)=P_{S/I}(t)-P_{S/I}(t-e).
\]
The leading \(t^r\) terms cancel, and the next term is
\[
a_rre\,t^{r-1}.
\]
Multiplying the leading coefficient by \((r-1)!\) gives
\[
e(r!a_r)=e\,\deg S/I.
\]
\end{solution}


\begin{problem}[CP-V-0086]
\label{prob:cp-v-0086}
Let
\[
I=(f_1,\dots,f_m)\subseteq k[x_1,\dots,x_n]
\]
be contained in
\[
\mathfrak m=(x_1,\dots,x_n),
\]
and put
\[
A=\bigl(k[x_1,\dots,x_n]/I\bigr)_{\bar{\mathfrak m}}.
\]
Let \(d=\dim A\). Prove that \(A\) is regular if and only if the Jacobian
matrix
\[
\left(
\frac{\partial f_i}{\partial x_j}(0)
\right)
\]
has rank
\[
n-d.
\]
\end{problem}

\begin{solution}
Let
\[
R=k[x_1,\dots,x_n]_{\mathfrak m}.
\]
Then
\[
A=R/IR
\]
and
\[
\mathfrak m_A/\mathfrak m_A^2
\cong
\mathfrak m/(\mathfrak m^2+I).
\]
The vector space
\[
\mathfrak m/\mathfrak m^2
\]
has basis given by the classes of \(x_1,\dots,x_n\). Modulo
\(\mathfrak m^2\), each \(f_i\) is its linear part:
\[
f_i\equiv
\sum_j
\frac{\partial f_i}{\partial x_j}(0)x_j
\pmod{\mathfrak m^2}.
\]
Hence the image of \(I\) in \(\mathfrak m/\mathfrak m^2\) has dimension
equal to the rank of the Jacobian matrix. Therefore
\[
\dim_k\mathfrak m_A/\mathfrak m_A^2
=
n-\operatorname{rank}J(0).
\]
By the cotangent-space characterization of regular local rings,
\[
A\text{ regular}
\iff
\dim_k\mathfrak m_A/\mathfrak m_A^2=d.
\]
Thus
\[
A\text{ regular}
\iff
\operatorname{rank}J(0)=n-d.
\]
\end{solution}



\section*{V/16 --- Support}

\begin{problem}[CP-V-0028]
\label{prob:cp-v-0028}
Let \(A\) be a ring and \(M\) an \(A\)-module. Recall
\[
\operatorname{Supp}(M)
=
\{\mathfrak p\in\operatorname{Spec}A:M_{\mathfrak p}\ne0\}.
\]

\textbf{(a)} If
\[
0\to M_1\to M_2\to M_3\to0
\]
is exact, prove that
\[
\operatorname{Supp}(M_2)
=
\operatorname{Supp}(M_1)\cup\operatorname{Supp}(M_3).
\]

\textbf{(b)} If \(M=\sum_i M_i\), prove that
\[
\operatorname{Supp}(M)
=
\bigcup_i\operatorname{Supp}(M_i).
\]

\textbf{(c)} If \(M\) is finitely generated, prove that
\[
\mathfrak p\in\operatorname{Supp}(M)
\iff
M\otimes_A k(\mathfrak p)\ne0.
\]

\textbf{(d)} If \(M,N\) are finitely generated, prove that
\[
\operatorname{Supp}(M\otimes_A N)
=
\operatorname{Supp}(M)\cap\operatorname{Supp}(N).
\]

\textbf{(e)} If \(f:A\to B\) is a ring homomorphism and \(M\) is finitely
generated, prove that
\[
\operatorname{Supp}_B(B\otimes_A M)
=
(f^*)^{-1}(\operatorname{Supp}_A M).
\]
\end{problem}

\begin{solution}
For \textbf{(a)}, localizing the short exact sequence at \(\mathfrak p\) gives
\[
0\to(M_1)_{\mathfrak p}
\to(M_2)_{\mathfrak p}
\to(M_3)_{\mathfrak p}
\to0.
\]
Thus \((M_2)_{\mathfrak p}=0\) exactly when both end terms vanish.

For \textbf{(b)}, the natural map
\[
\bigoplus_iM_i\to M
\]
is surjective. After localizing at \(\mathfrak p\), the target is nonzero
exactly when at least one localized summand is nonzero.

For \textbf{(c)}, put \(N=M_{\mathfrak p}\). Then
\[
M\otimes_Ak(\mathfrak p)
\cong
N/\mathfrak pA_{\mathfrak p}N.
\]
Since \(N\) is finitely generated over the local ring \(A_{\mathfrak p}\),
Nakayama's lemma gives
\[
N=0
\iff
N/\mathfrak pA_{\mathfrak p}N=0.
\]

For \textbf{(d)},
\[
(M\otimes_AN)_{\mathfrak p}
\cong
M_{\mathfrak p}\otimes_{A_{\mathfrak p}}N_{\mathfrak p}.
\]
If either factor vanishes, so does the tensor product. Conversely, if both
are nonzero, then by part \textbf{(c)} their fibers over \(k(\mathfrak p)\)
are nonzero finite-dimensional vector spaces. Their tensor product over the
field \(k(\mathfrak p)\) is nonzero, so the localized tensor product is
nonzero.

For \textbf{(e)}, let \(\mathfrak q\in\operatorname{Spec}B\) and put
\[
\mathfrak p=f^{-1}(\mathfrak q).
\]
Then
\[
(B\otimes_AM)\otimes_Bk(\mathfrak q)
\cong
M\otimes_Ak(\mathfrak q)
\cong
(M\otimes_Ak(\mathfrak p))
\otimes_{k(\mathfrak p)}k(\mathfrak q).
\]
Since extension of scalars along a field extension preserves nonzero vector
spaces, the last module is nonzero exactly when
\[
M\otimes_Ak(\mathfrak p)\ne0.
\]
Part \textbf{(c)} completes the proof.
\end{solution}


\begin{problem}[CP-V-0107]
\label{prob:cp-v-0107}
Let \(A\) be a ring and \(M\) a finitely generated \(A\)-module. Prove that
\[
\operatorname{Supp}(M)=V(\operatorname{Ann}_A M).
\]
Deduce that \(\operatorname{Supp}(M)\) is closed in \(\operatorname{Spec}A\).
\end{problem}

\begin{solution}
If \(\mathfrak p\not\supseteq\operatorname{Ann}(M)\), choose
\[
a\in\operatorname{Ann}(M)\setminus\mathfrak p.
\]
Then \(a/1\) is a unit in \(A_{\mathfrak p}\), while it annihilates \(M_{\mathfrak p}\). Hence
\[
M_{\mathfrak p}=0.
\]

Conversely, assume \(M_{\mathfrak p}=0\). Choose generators \(m_1,\dots,m_r\) of \(M\). For each \(i\), some
\[
s_i\notin\mathfrak p
\]
satisfies \(s_im_i=0\). Put \(s=s_1\cdots s_r\). Then \(s\notin\mathfrak p\) and \(sM=0\), so
\[
s\in\operatorname{Ann}(M)\setminus\mathfrak p.
\]
Thus \(\mathfrak p\notin V(\operatorname{Ann}M)\). Therefore
\[
\operatorname{Supp}(M)=V(\operatorname{Ann}M),
\]
which is closed.
\end{solution}



\section*{V/17 --- Associated Primes}

\begin{problem}[CP-V-0026]
\label{prob:cp-v-0026}
Let \(X\) be a Noetherian topological space.

\textbf{(a)} Prove that every closed subset \(Z\subseteq X\) can be written as a
finite union
\[
Z=Z_1\cup\cdots\cup Z_r
\]
of irreducible closed subsets, and that an irredundant such decomposition is
unique up to reordering.

\textbf{(b)} Let \(A\) be a ring. Prove that if \(\mathfrak p\subseteq A\) is
prime, then
\[
V(\mathfrak p)\subseteq\operatorname{Spec}A
\]
is irreducible.

\textbf{(c)} Prove that if \(A\) is Noetherian, then
\[
\operatorname{Spec}A
\]
is a Noetherian topological space.

\textbf{(d)} Assume \(A\) is Noetherian and \(I\subseteq A\). Show that the
irreducible components of \(V(I)\) correspond bijectively to the minimal
primes over \(I\), equivalently to the minimal associated primes of \(A/I\).
\end{problem}

\begin{solution}
For \textbf{(a)}, if \(Z\) is irreducible there is nothing to prove. If it is
reducible, write it as a union of two proper closed subsets and continue.
An infinite continuation would produce an infinite strictly descending chain
of closed subsets, impossible in a Noetherian space. Hence the process stops
after finitely many steps.

For uniqueness, suppose
\[
Z=\bigcup_{i=1}^r Z_i=\bigcup_{j=1}^s W_j
\]
are irredundant decompositions into irreducible closed subsets. Fix \(i\).
Since
\[
Z_i=\bigcup_j(Z_i\cap W_j),
\]
irreducibility implies \(Z_i\subseteq W_j\) for some \(j\). Reversing the
roles gives \(W_j\subseteq Z_k\) for some \(k\), hence
\[
Z_i\subseteq W_j\subseteq Z_k.
\]
Irredundancy forces \(k=i\), so \(Z_i=W_j\). Symmetry gives equality of the
two collections.

For \textbf{(b)}, suppose
\[
V(\mathfrak p)=V(I)\cup V(J).
\]
Then
\[
V(\mathfrak p)=V(IJ),
\]
so
\[
\sqrt{IJ}=\mathfrak p.
\]
Thus \(IJ\subseteq\mathfrak p\). Since \(\mathfrak p\) is prime,
\(I\subseteq\mathfrak p\) or \(J\subseteq\mathfrak p\), and hence one of the
two closed subsets equals all of \(V(\mathfrak p)\). Therefore
\(V(\mathfrak p)\) is irreducible.

For \textbf{(c)}, a descending chain
\[
V(I_1)\supseteq V(I_2)\supseteq\cdots
\]
induces an ascending chain
\[
\sqrt{I_1}\subseteq\sqrt{I_2}\subseteq\cdots.
\]
Because \(A\) is Noetherian, this chain stabilizes, and therefore the closed
sets stabilize.

For \textbf{(d)}, every irreducible closed subset of \(\operatorname{Spec}A\)
has the form \(V(\mathfrak p)\) for a prime \(\mathfrak p\). Maximal
irreducible closed subsets of \(V(I)\) therefore correspond exactly to primes
\(\mathfrak p\supseteq I\) minimal with respect to inclusion. These are the
minimal primes over \(I\). In a Noetherian ring the minimal primes over \(I\)
are exactly the minimal elements of
\[
\operatorname{Ass}(A/I).
\]
Hence the irreducible components of \(V(I)\) correspond precisely to the
minimal associated primes of \(A/I\).
\end{solution}


\begin{problem}[CP-V-0027]
\label{prob:cp-v-0027}
Let
\[
R=k[x,y,z],
\qquad
I=(x^2,xy,xz,yz).
\]
Determine
\[
\operatorname{Ass}(R/I)
\]
and find an irredundant primary decomposition of \(I\).
\end{problem}

\begin{solution}
If a prime \(\mathfrak p\) contains \(I\), then \(x^2\in\mathfrak p\), hence
\(x\in\mathfrak p\). Since \(yz\in\mathfrak p\), primality implies
\(y\in\mathfrak p\) or \(z\in\mathfrak p\). Thus the minimal primes over
\(I\) are
\[
(x,y),\qquad(x,z).
\]

Now consider the class \(\bar x\in R/I\). Since
\[
x^2,xy,xz\in I,
\]
every element of \((x,y,z)\) annihilates \(\bar x\). Conversely, if
\(fx\in I\), inspection of monomials shows that \(f\in(x,y,z)\). Hence
\[
\operatorname{Ann}_{R/I}(\bar x)=(x,y,z),
\]
so the maximal ideal \((x,y,z)\) is an embedded associated prime. Therefore
\[
\operatorname{Ass}(R/I)
=
\{(x,y),(x,z),(x,y,z)\}.
\]

For the decomposition,
\[
(x,y)\cap(x,z)=(x,yz).
\]
Intersecting with \((x^2,y,z)\) gives the monomial ideal generated by
\[
x^2,\ xy,\ xz,\ yz.
\]
Hence
\[
I=(x,y)\cap(x,z)\cap(x^2,y,z).
\]
The first two ideals are prime, hence primary. Also
\[
R/(x^2,y,z)\cong k[x]/(x^2),
\]
so \((x^2,y,z)\) is \((x,y,z)\)-primary. The decomposition is irredundant,
and therefore
\[
\boxed{
I=(x,y)\cap(x,z)\cap(x^2,y,z)
}.
\]
\end{solution}


\begin{problem}[CP-V-0029]
\label{prob:cp-v-0029}
Let \(A\) be a ring and let \(\mathfrak q\subseteq A\) be
\(\mathfrak p\)-primary.

\textbf{(a)} Prove that
\[
\mathfrak q[x]\subseteq A[x]
\]
is \(\mathfrak p[x]\)-primary.

\textbf{(b)} If
\[
I=\bigcap_{i=1}^n\mathfrak q_i
\]
is an irredundant primary decomposition in \(A\), prove that
\[
I[x]
=
\bigcap_{i=1}^n\mathfrak q_i[x]
\]
is an irredundant primary decomposition in \(A[x]\).

\textbf{(c)} If \(\mathfrak p\) is minimal over \(I\), prove that
\[
\mathfrak p[x]
\]
is minimal over \(I[x]\).
\end{problem}

\begin{solution}
We have
\[
A[x]/\mathfrak q[x]\cong(A/\mathfrak q)[x]
\]
and
\[
\sqrt{\mathfrak q[x]}
=
\sqrt{\mathfrak q}[x]
=
\mathfrak p[x].
\]
Since \(\mathfrak q\) is primary, every zero-divisor of \(A/\mathfrak q\) is
nilpotent. If
\[
f(x)\in(A/\mathfrak q)[x]
\]
is a zero-divisor, McCoy's theorem gives a nonzero \(a\in A/\mathfrak q\)
with \(af(x)=0\). Hence every coefficient of \(f\) is a zero-divisor in
\(A/\mathfrak q\), and therefore nilpotent. A polynomial with finitely many
nilpotent coefficients is nilpotent. Thus every zero-divisor of
\((A/\mathfrak q)[x]\) is nilpotent, so \(\mathfrak q[x]\) is
\(\mathfrak p[x]\)-primary.

For \textbf{(b)}, coefficientwise membership gives
\[
\left(\bigcap_i\mathfrak q_i\right)[x]
=
\bigcap_i\mathfrak q_i[x].
\]
Part \textbf{(a)} shows that each extended component remains primary.
If one component became redundant, say
\[
\bigcap_{i\ne j}\mathfrak q_i[x]\subseteq\mathfrak q_j[x],
\]
then taking a constant polynomial
\[
a\in\bigcap_{i\ne j}\mathfrak q_i\setminus\mathfrak q_j
\]
would give a contradiction. Hence irredundancy is preserved.

For \textbf{(c)}, \(\mathfrak p[x]\) is prime and contains \(I[x]\).
Suppose
\[
I[x]\subseteq P\subseteq\mathfrak p[x]
\]
with \(P\) prime. Contracting to \(A\) gives
\[
I\subseteq P\cap A\subseteq\mathfrak p.
\]
By minimality of \(\mathfrak p\),
\[
P\cap A=\mathfrak p.
\]
Hence \(\mathfrak p[x]\subseteq P\). Since also \(P\subseteq\mathfrak p[x]\),
we obtain
\[
P=\mathfrak p[x].
\]
Therefore \(\mathfrak p[x]\) is minimal over \(I[x]\).
\end{solution}


\begin{problem}[CP-V-0051]
\label{prob:cp-v-0051}
Let \(A\) be a ring and let \(D(A)\) be the set of prime ideals
\(\mathfrak p\) that are minimal over \(\operatorname{Ann}(a)\) for some
\(a\in A\).

\textbf{(a)} Prove that \(x\in A\) is a zero divisor if and only if
\[
x\in\mathfrak p
\]
for some \(\mathfrak p\in D(A)\).

\textbf{(b)} If \(S\subseteq A\) is multiplicatively closed, prove
\[
D(S^{-1}A)=D(A)\cap\operatorname{Spec}(S^{-1}A).
\]

\textbf{(c)} If the zero ideal has a primary decomposition, prove that
\[
D(A)=\operatorname{Ass}_A(A).
\]
\end{problem}

\begin{solution}
If \(x\) is a zero divisor, choose \(0\ne a\) with \(xa=0\). Then
\(x\in\operatorname{Ann}(a)\), hence \(x\) lies in a prime minimal over
\(\operatorname{Ann}(a)\).

Conversely, suppose \(x\in\mathfrak p\in D(A)\), with \(\mathfrak p\)
minimal over \(I=\operatorname{Ann}(a)\). In \(B=A/I\), the image of
\(\mathfrak p\) is a minimal prime. Localizing at it makes every element of
that prime nilpotent. Thus some nonzero class is annihilated by \(\bar x\).
Lifting back to \(A\) produces \(0\ne c\in A\) with \(xc=0\). Hence \(x\) is
a zero divisor.

For localization, annihilators localize:
\[
\operatorname{Ann}_{S^{-1}A}(a/1)
=
S^{-1}\operatorname{Ann}_A(a),
\]
and minimal primes disjoint from \(S\) correspond bijectively under
extension and contraction. This gives
\[
D(S^{-1}A)=D(A)\cap\operatorname{Spec}(S^{-1}A).
\]

If
\[
0=Q_1\cap\cdots\cap Q_n
\]
is irredundant primary with radicals \(\mathfrak p_i\), then for each \(i\)
one can choose
\[
a_i\in\bigcap_{j\ne i}Q_j\setminus Q_i
\]
so that
\[
\sqrt{\operatorname{Ann}(a_i)}=\mathfrak p_i.
\]
Hence every associated prime lies in \(D(A)\). Conversely, any prime minimal
over an annihilator arises among the radicals of the relevant colon-primary
components, so it is one of the \(\mathfrak p_i\). Therefore
\[
D(A)=\operatorname{Ass}_A(A).
\]
\end{solution}



\section*{V/18 --- Completion and I-Adic Topology}

\begin{problem}[CP-V-0030]
\label{prob:cp-v-0030}
Let \(A\) be a ring and \(I\subseteq A\) an ideal. Define the \(I\)-adic
completion
\[
\widehat A=\varprojlim_n A/I^n.
\]

\textbf{(a)} Show that the \((x_1,\dots,x_r)\)-adic completion of
\[
k[x_1,\dots,x_r]
\]
is
\[
k[[x_1,\dots,x_r]].
\]

\textbf{(b)} Show that the \((p)\)-adic completion of \(\mathbb Z\) is
\[
\mathbb Z_p=\varprojlim_n\mathbb Z/p^n\mathbb Z.
\]

\textbf{(c)} Let \(M\) be an \(A\)-module. Show that the canonical map
\[
M\longrightarrow\widehat M:=\varprojlim_n M/I^nM
\]
has kernel
\[
\bigcap_{n\ge0}I^nM.
\]
\end{problem}

\begin{solution}
For \textbf{(a)}, modulo
\[
\mathfrak m^n,\qquad \mathfrak m=(x_1,\dots,x_r),
\]
a polynomial is determined by its terms of total degree \(<n\). A compatible
family
\[
(f_n+\mathfrak m^n)_n
\]
therefore determines, and is determined by, one coefficient for every
monomial. Thus
\[
\varprojlim_n k[x_1,\dots,x_r]/\mathfrak m^n
\cong
k[[x_1,\dots,x_r]].
\]

For \textbf{(b)}, an element of
\[
\varprojlim_n\mathbb Z/p^n\mathbb Z
\]
is a compatible family
\[
(a_n)_n,\qquad
a_{n+1}\equiv a_n\pmod{p^n}.
\]
Equivalently it is represented uniquely by a \(p\)-adic expansion
\[
a_0+a_1p+a_2p^2+\cdots,
\qquad
0\le a_i<p.
\]
This is precisely the ring \(\mathbb Z_p\).

For \textbf{(c)}, define
\[
\iota_M(m)=\bigl(m+I^nM\bigr)_n.
\]
Then \(\iota_M(m)=0\) exactly when
\[
m\in I^nM
\]
for every \(n\). Hence
\[
\ker(\iota_M)=\bigcap_{n\ge0}I^nM.
\]
\end{solution}


\begin{problem}[CP-V-0031]
\label{prob:cp-v-0031}
Let \(A\) be Noetherian, \(I\subseteq A\) an ideal, and
\[
0\longrightarrow M_1\longrightarrow M_2\longrightarrow M_3\longrightarrow0
\]
a short exact sequence of finitely generated \(A\)-modules.

\textbf{(a)} Prove that \(I\)-adic completion preserves exactness:
\[
0\longrightarrow\widehat M_1
\longrightarrow\widehat M_2
\longrightarrow\widehat M_3
\longrightarrow0.
\]

\textbf{(b)} Prove that for every finitely generated \(A\)-module \(M\),
the natural map
\[
\widehat A\otimes_A M\longrightarrow\widehat M
\]
is an isomorphism.
\end{problem}

\begin{solution}
For each \(n\), the original short exact sequence induces
\[
0\longrightarrow
\frac{M_1}{M_1\cap I^nM_2}
\longrightarrow
\frac{M_2}{I^nM_2}
\longrightarrow
\frac{M_3}{I^nM_3}
\longrightarrow0.
\]
The transition maps in the left-hand inverse system are surjective, so
passing to inverse limits preserves exactness on the right. Hence
\[
0\longrightarrow
\varprojlim_n\frac{M_1}{M_1\cap I^nM_2}
\longrightarrow
\widehat M_2
\longrightarrow
\widehat M_3
\longrightarrow0.
\]

By the Artin--Rees lemma, the filtrations
\[
\{I^nM_1\}
\qquad\text{and}\qquad
\{M_1\cap I^nM_2\}
\]
are equivalent. Therefore they define the same completion:
\[
\varprojlim_n\frac{M_1}{M_1\cap I^nM_2}
\cong
\widehat M_1.
\]
This proves \textbf{(a)}.

For \textbf{(b)}, choose a finite presentation
\[
A^r\longrightarrow A^s\longrightarrow M\longrightarrow0.
\]
Completion is exact by part \textbf{(a)}, and
\[
\widehat{A^m}\cong\widehat A^{\,m}
\cong
\widehat A\otimes_A A^m.
\]
Thus both \(\widehat M\) and \(\widehat A\otimes_A M\) are cokernels of the
same map
\[
\widehat A^{\,r}\longrightarrow\widehat A^{\,s}.
\]
Hence
\[
\widehat A\otimes_A M\cong\widehat M.
\]
\end{solution}


\begin{problem}[CP-V-0032]
\label{prob:cp-v-0032}
Let \(A\) be Noetherian, \(I\subseteq A\), and let \(\widehat A\) be the
\(I\)-adic completion. Denote by
\[
\widehat I
\]
the completion of \(I\) as an \(A\)-module.

\textbf{(a)} Prove that
\[
\widehat I=I\widehat A
\cong
\widehat A\otimes_A I.
\]

\textbf{(b)} Prove that
\[
\widehat{I^n}=(\widehat I)^n.
\]

\textbf{(c)} Prove that
\[
I^n/I^{n+1}
\cong
\widehat I^{\,n}/\widehat I^{\,n+1}.
\]

\textbf{(d)} Show that
\[
\widehat I\subseteq\operatorname{Jac}(\widehat A).
\]

\textbf{(e)} Deduce that if \((A,\mathfrak m)\) is Noetherian local and
\(\widehat A\) is the \(\mathfrak m\)-adic completion, then \(\widehat A\)
is local with maximal ideal
\[
\mathfrak m\widehat A.
\]
\end{problem}

\begin{solution}
Apply exactness of completion to
\[
0\longrightarrow I\longrightarrow A\longrightarrow A/I\longrightarrow0.
\]
Since \(I(A/I)=0\), the module \(A/I\) is already \(I\)-adically complete.
Thus
\[
0\longrightarrow\widehat I
\longrightarrow\widehat A
\longrightarrow A/I
\longrightarrow0.
\]
The kernel of \(\widehat A\to A/I\) is the extended ideal \(I\widehat A\),
so
\[
\widehat I=I\widehat A.
\]
Also, by the tensor description of completion,
\[
\widehat I\cong\widehat A\otimes_A I.
\]

Applying the same argument to \(I^n\) gives
\[
\widehat{I^n}
=
I^n\widehat A
=
(I\widehat A)^n
=
(\widehat I)^n.
\]

Now complete
\[
0\longrightarrow I^{n+1}
\longrightarrow I^n
\longrightarrow I^n/I^{n+1}
\longrightarrow0.
\]
Because \(I\) annihilates \(I^n/I^{n+1}\), that quotient is already complete.
Therefore
\[
0\longrightarrow
\widehat I^{\,n+1}
\longrightarrow
\widehat I^{\,n}
\longrightarrow
I^n/I^{n+1}
\longrightarrow0,
\]
which gives
\[
I^n/I^{n+1}
\cong
\widehat I^{\,n}/\widehat I^{\,n+1}.
\]

If \(u\in\widehat I\), then the geometric series
\[
1+u+u^2+\cdots
\]
converges in the \(I\)-adic topology of the complete ring \(\widehat A\).
Its sum is an inverse to \(1-u\). Hence every \(1-u\) with
\(u\in\widehat I\) is a unit, so
\[
\widehat I\subseteq\operatorname{Jac}(\widehat A).
\]

Finally, take \(I=\mathfrak m\). We have
\[
\widehat A/\mathfrak m\widehat A
\cong
A/\mathfrak m,
\]
which is a field. Thus \(\mathfrak m\widehat A\) is maximal. Since it lies
in the Jacobson radical, every maximal ideal of \(\widehat A\) contains it.
Therefore it is the unique maximal ideal, and \(\widehat A\) is local.
\end{solution}


\begin{problem}[CP-V-0053]
\label{prob:cp-v-0053}
Let \(A\) be Noetherian and \(I\subseteq A\). Define
\[
\operatorname{gr}_I(A)
=
\bigoplus_{n\ge0}I^n/I^{n+1}.
\]

\textbf{(a)} Prove that \(\operatorname{gr}_I(A)\) is Noetherian.

\textbf{(b)} If \((A,\mathfrak m)\) is Noetherian local and
\(\operatorname{gr}_{\mathfrak m}(A)\) is an integral domain, prove that
\(A\) is an integral domain.
\end{problem}

\begin{solution}
Write
\[
I=(a_1,\dots,a_r).
\]
The initial classes
\[
\overline a_i\in I/I^2
\]
generate the associated graded ring over \(A/I\), giving a surjection
\[
(A/I)[T_1,\dots,T_r]\twoheadrightarrow\operatorname{gr}_I(A).
\]
Hence \(\operatorname{gr}_I(A)\) is Noetherian.

For \textbf{(b)}, let \(0\ne a,b\in A\). By Krull intersection, choose
largest integers \(r,s\) such that
\[
a\in\mathfrak m^r\setminus\mathfrak m^{r+1},
\qquad
b\in\mathfrak m^s\setminus\mathfrak m^{s+1}.
\]
Their initial forms
\[
a^*\in\mathfrak m^r/\mathfrak m^{r+1},
\qquad
b^*\in\mathfrak m^s/\mathfrak m^{s+1}
\]
are nonzero. Since the associated graded ring is a domain,
\[
a^*b^*\ne0.
\]
Thus
\[
ab\notin\mathfrak m^{r+s+1},
\]
in particular \(ab\ne0\). Hence \(A\) is a domain.
\end{solution}


\begin{problem}[CP-V-0054]
\label{prob:cp-v-0054}
Let
\[
G\supseteq G_1\supseteq G_2\supseteq\cdots
\]
be a descending filtration of an abelian group \(G\), and give \(G\) the
linear topology whose basic neighborhoods of \(0\) are the \(G_n\).

\textbf{(a)} Prove that this topology is Hausdorff if
\[
\bigcap_{n\ge1}G_n=0.
\]

\textbf{(b)} Define the completion \(\widehat G\) by Cauchy sequences modulo
equivalence. Prove that
\[
\ker(G\to\widehat G)=\bigcap_{n\ge1}G_n.
\]

\textbf{(c)} Prove the natural isomorphism
\[
\widehat G\cong\varprojlim_n G/G_n.
\]
\end{problem}

\begin{solution}
If \(x\ne y\), then \(x-y\notin G_n\) for some \(n\), and the cosets
\[
x+G_n,\qquad y+G_n
\]
are disjoint. Hence the topology is Hausdorff.

The constant sequence attached to \(g\in G\) represents zero in the
completion exactly when for every \(n\),
\[
g\in G_n.
\]
Thus the kernel is
\[
\bigcap_nG_n.
\]

A Cauchy sequence \((x_r)\) is eventually constant modulo each \(G_n\), so
it determines a compatible family
\[
(\bar x_n)_n\in\varprojlim G/G_n.
\]
Equivalent Cauchy sequences determine the same family.

Conversely, from a compatible family \((\bar x_n)\), choose representatives
\(x_n\in G\). Compatibility implies
\[
x_{n+1}-x_n\in G_n,
\]
so \((x_n)\) is Cauchy. Different representative choices give equivalent
Cauchy sequences. The two constructions are inverse.
\end{solution}


\begin{problem}[CP-V-0055]
\label{prob:cp-v-0055}
Assume \(\operatorname{char}k\ne2\). Prove that
\[
A=
k[x,y]/\bigl(y^2-x^2(1+x)\bigr)
\]
is an integral domain, but its \((x,y)\)-adic completion
\[
\widehat A
\cong
k[[x,y]]/\bigl(y^2-x^2(1+x)\bigr)
\]
is not.
\end{problem}

\begin{solution}
Set
\[
f=y^2-x^2(1+x).
\]
Viewed as a quadratic polynomial in \(y\) over \(k(x)\), \(f\) is reducible
only if \(1+x\) is a square in \(k(x)\). But the valuation at \(x+1\) of
\(1+x\) is \(1\), whereas every square has even valuation. Thus \(f\) is
irreducible in \(k(x)[y]\), hence by Gauss's lemma in \(k[x,y]\). Therefore
\((f)\) is prime and \(A\) is a domain.

In \(k[[x]]\), the unit \(1+x\) has a square root
\[
u=(1+x)^{1/2}
\]
because \(\operatorname{char}k\ne2\). Hence
\[
f=(y-xu)(y+xu)
\]
in \(k[[x,y]]\). Neither factor lies in \((f)\), but their product does.
Thus the completed quotient has nonzero zero divisors and is not a domain.
\end{solution}


\begin{problem}[CP-V-0087]
\label{prob:cp-v-0087}
Let
\[
0\to A_n\to B_n\to C_n\to0
\]
be a compatible short exact sequence of inverse systems of abelian groups.

\textbf{(a)} Prove that inverse limit is left exact:
\[
0\to\varprojlim A_n
\to\varprojlim B_n
\to\varprojlim C_n.
\]

\textbf{(b)} If every transition map
\[
A_{n+1}\to A_n
\]
is surjective, prove that the right-hand map is also surjective.
\end{problem}

\begin{solution}
The induced maps on inverse limits are coordinatewise. Injectivity on the
left follows from injectivity of every \(A_n\to B_n\).

If \((b_n)\) maps to zero in \(\varprojlim C_n\), then each \(b_n\) lies in
the image of a unique \(a_n\in A_n\). Compatibility of the \(b_n\), together
with injectivity of \(A_n\to B_n\), forces the \(a_n\) to be compatible.
Hence
\[
\ker(\varprojlim B_n\to\varprojlim C_n)
=
\operatorname{im}(\varprojlim A_n).
\]

For surjectivity, take a compatible \((c_n)\). Choose \(b_1\) lifting
\(c_1\). Inductively, choose any lift \(b'_{n+1}\) of \(c_{n+1}\). The error
between its image in \(B_n\) and the already chosen \(b_n\) lies in the image
of \(A_n\). Lift that correction to \(A_{n+1}\) using surjectivity of
\(A_{n+1}\to A_n\), subtract it from \(b'_{n+1}\), and obtain a compatible
lift \(b_{n+1}\). Thus every compatible \((c_n)\) lifts.
\end{solution}


\begin{problem}[CP-V-0088]
\label{prob:cp-v-0088}
Let \((A,\mathfrak m)\) be local and let \(\widehat A\) be its
\(\mathfrak m\)-adic completion. Let
\[
f(T)\in\widehat A[T]
\]
be monic. Suppose \(\bar a\in A/\mathfrak m\) satisfies
\[
\bar f(\bar a)=0,
\qquad
\bar f'(\bar a)\ne0.
\]
Prove that there exists
\[
a\in\widehat A
\]
with
\[
f(a)=0.
\]
\end{problem}

\begin{solution}
Choose a lift \(a_1\in\widehat A\) of \(\bar a\). Then
\[
f(a_1)\in\widehat{\mathfrak m},
\]
while
\[
f'(a_1)\notin\widehat{\mathfrak m},
\]
so \(f'(a_1)\) is a unit.

Inductively suppose
\[
f(a_r)\in\widehat{\mathfrak m}^{\,r}.
\]
Seek
\[
a_{r+1}=a_r+t_r,
\qquad
t_r\in\widehat{\mathfrak m}^{\,r}.
\]
Modulo \(\widehat{\mathfrak m}^{\,r+1}\),
\[
f(a_r+t_r)
\equiv
f(a_r)+t_rf'(a_r),
\]
because \(t_r^2\in\widehat{\mathfrak m}^{\,2r}\subseteq
\widehat{\mathfrak m}^{\,r+1}\). Since \(f'(a_r)\) remains a unit, choose
\(t_r\) so that
\[
f(a_{r+1})\in\widehat{\mathfrak m}^{\,r+1}.
\]

Then
\[
a_{r+1}-a_r\in\widehat{\mathfrak m}^{\,r},
\]
so \((a_r)\) is Cauchy. Completeness gives \(a_r\to a\in\widehat A\).
Continuity of polynomial evaluation yields
\[
f(a)=\lim_r f(a_r)=0.
\]
\end{solution}


\begin{problem}[CP-V-0089]
\label{prob:cp-v-0089}
Let
\[
A=k[[x,y]]/(x^2-y^3),
\qquad
\mathfrak m=(x,y).
\]
Prove that
\[
\operatorname{gr}_{\mathfrak m}(A)
\]
is not an integral domain.
\end{problem}

\begin{solution}
Let \(x^*,y^*\) be the initial forms of \(x,y\) in
\[
\mathfrak m/\mathfrak m^2.
\]
Both are nonzero. But in \(A\),
\[
x^2=y^3\in\mathfrak m^3.
\]
Thus the class of \(x^2\) in
\[
\mathfrak m^2/\mathfrak m^3
\]
vanishes. Equivalently,
\[
(x^*)^2=0
\]
in the associated graded ring. Hence \(x^*\) is a nonzero nilpotent, so
\[
\operatorname{gr}_{\mathfrak m}(A)
\]
is not a domain.
\end{solution}



\section*{V/19 --- Integral Dependence}

\begin{problem}[CP-V-0108]
\label{prob:cp-v-0108}
Let \(A\subseteq B\) be rings. Prove that the set of elements of \(B\) integral over \(A\) is a subring of \(B\).
\end{problem}

\begin{solution}
Let \(x,y\in B\) be integral over \(A\). Then \(A[x]\) and \(A[y]\) are finite \(A\)-modules. Hence
\[
A[x,y]
\]
is finite over \(A[x]\), and therefore finite over \(A\). Both
\[
x+y,\qquad xy
\]
belong to the finite \(A\)-algebra \(A[x,y]\). Every element of a finite \(A\)-algebra is integral over \(A\), so \(x+y\) and \(xy\) are integral. Also \(-x\) is integral. Thus the integral elements form a subring.
\end{solution}


\begin{problem}[CP-V-0109]
\label{prob:cp-v-0109}
Let \(A\subseteq B\subseteq C\). If \(B\) is integral over \(A\) and \(C\) is integral over \(B\), prove that \(C\) is integral over \(A\).
\end{problem}

\begin{solution}
Take \(c\in C\). Since \(c\) is integral over \(B\), it satisfies a monic equation
\[
c^n+b_{n-1}c^{n-1}+\cdots+b_0=0
\]
with \(b_i\in B\). Let
\[
B_0=A[b_0,\dots,b_{n-1}].
\]
Because each \(b_i\) is integral over \(A\), \(B_0\) is finite as an \(A\)-module. The relation shows that
\[
B_0[c]
\]
is finite as a \(B_0\)-module, hence finite as an \(A\)-module. Therefore \(c\) is integral over \(A\).
\end{solution}



\section*{V/20 --- Integral Closure and Normalization}

\begin{problem}[CP-V-0033]
\label{prob:cp-v-0033}
Let \(A\) be a unique factorization domain with fraction field \(K\).

\textbf{(a)} Prove that \(A\) is integrally closed in \(K\).

\textbf{(b)} Let
\[
A=k[x,y]/(y^2-x^3).
\]
Show that \(A\) is not normal.

\textbf{(c)} Identify the integral closure of \(A\) in its fraction field.
\end{problem}

\begin{solution}
For \textbf{(a)}, let
\[
z=\frac ab\in K
\]
with \(a,b\in A\) coprime and suppose \(z\) is integral over \(A\). Then
\[
z^n+c_{n-1}z^{n-1}+\cdots+c_0=0
\]
with \(c_i\in A\). Multiplying by \(b^n\) gives
\[
a^n
=
-b\bigl(c_{n-1}a^{n-1}+c_{n-2}a^{n-2}b+\cdots+c_0b^{n-1}\bigr).
\]
Hence \(b\mid a^n\). Since \(a\) and \(b\) are coprime in a UFD, \(b\) is a
unit. Thus \(z\in A\), so \(A\) is integrally closed.

For \textbf{(b)}, write
\[
A\cong k[t^2,t^3]\subseteq k[t]
\]
via
\[
x\mapsto t^2,\qquad y\mapsto t^3.
\]
Then
\[
t=\frac{t^3}{t^2}
\]
lies in \(\operatorname{Frac}(A)\), and it is integral over \(A\) because it
satisfies
\[
T^2-x=0.
\]
But \(t\notin k[t^2,t^3]\). Hence \(A\) is not integrally closed.

For \textbf{(c)}, the ring \(k[t]\) is a PID, hence a UFD, hence integrally
closed. Since \(t\) is integral over \(A\), the extension
\[
A\subseteq k[t]
\]
is integral. Therefore the integral closure of \(A\) in its fraction field
\(k(t)\) is
\[
\boxed{k[t]}.
\]
\end{solution}


\begin{problem}[CP-V-0034]
\label{prob:cp-v-0034}
Let \(A\subseteq B\) be rings and suppose that
\[
B\setminus A
\]
is multiplicatively closed. Prove that \(A\) is integrally closed in \(B\);
that is, every element of \(B\) integral over \(A\) already belongs to \(A\).
\end{problem}

\begin{solution}
Suppose \(b\in B\) is integral over \(A\), and assume for contradiction that
\(b\notin A\). Then \(b\in B\setminus A\). Since the complement is
multiplicatively closed,
\[
b^m\in B\setminus A
\]
for every \(m\ge1\).

On the other hand, integrality gives a monic relation
\[
b^n+a_{n-1}b^{n-1}+\cdots+a_1b+a_0=0,
\qquad a_i\in A.
\]
Hence
\[
b\bigl(b^{n-1}+a_{n-1}b^{n-2}+\cdots+a_1\bigr)=-a_0\in A.
\]
If the second factor were outside \(A\), then by multiplicative closure of
\(B\setminus A\) the product would lie outside \(A\), a contradiction.
Thus
\[
b^{n-1}+a_{n-1}b^{n-2}+\cdots+a_1\in A.
\]
Subtracting the lower terms shows
\[
b^{n-1}\in A.
\]
But this contradicts the fact that every positive power of \(b\) lies in
\(B\setminus A\). Therefore \(b\in A\).

So \(A\) is integrally closed in \(B\).
\end{solution}



\section*{V/21 --- Valuation Rings}

\begin{problem}[CP-V-0035]
\label{prob:cp-v-0035}
Let \(K\) be a field. Consider the set \(\Sigma\) of local subrings of \(K\),
ordered by domination: \(B\) dominates \(A\) if
\[
A\subseteq B
\]
and the maximal ideal of \(B\) contracts to the maximal ideal of \(A\).

\textbf{(a)} Prove that \(\Sigma\) has maximal elements.

\textbf{(b)} Prove that a maximal element \(V\in\Sigma\) is a valuation ring
of \(K\); that is, for every \(x\in K^\times\),
\[
x\in V
\quad\text{or}\quad
x^{-1}\in V.
\]

\textbf{(c)} Conversely, prove that every valuation ring \(V\subseteq K\) is
maximal in \(\Sigma\).
\end{problem}

\begin{solution}
For \textbf{(a)}, let
\[
A_1\subseteq A_2\subseteq\cdots
\]
be a chain in \(\Sigma\). Its union
\[
A=\bigcup_iA_i
\]
is a subring of \(K\). If \(\mathfrak m_i\) is the maximal ideal of \(A_i\),
then
\[
\mathfrak m=\bigcup_i\mathfrak m_i
\]
is an ideal of \(A\), and every element outside \(\mathfrak m\) is a unit in
some \(A_i\), hence in \(A\). Thus \(A\) is local and dominates every ring in
the chain. Zorn's lemma gives a maximal element.

For \textbf{(b)}, let \(V\) be maximal and suppose \(x\notin V\). If
\(x^{-1}\notin V\) as well, then \(V[x^{-1}]\) is integral over \(V\) after
using a relation arising from the maximality/domination argument in the
source construction. By lying over, one obtains a prime of \(V[x^{-1}]\)
lying over the maximal ideal of \(V\). Localizing at that prime gives a local
subring of \(K\) properly dominating \(V\), contradicting maximality.
Therefore
\[
x\notin V\implies x^{-1}\in V.
\]
Hence \(V\) is a valuation ring.

For \textbf{(c)}, let \(V\subseteq B\subseteq K\) with \(B\) local and
dominating \(V\). Take \(x\in B\). If \(x\notin V\), valuation-ring
comparability gives
\[
x^{-1}\in V.
\]
Since \(x\in B\), this makes \(x\) a unit of \(B\). But domination forces
elements of the maximal ideal of \(V\) to remain nonunits in \(B\), and the
valuation-ring dichotomy excludes a genuine enlargement preserving the same
center. Hence every \(x\in B\) lies in \(V\), so
\[
B=V.
\]
Thus valuation rings are precisely the maximal local subrings of \(K\) under
domination.
\end{solution}


\begin{problem}[CP-V-0036]
\label{prob:cp-v-0036}
Let \(K\) be a field, let \(\Gamma\) be a totally ordered abelian group, and
let
\[
v:K^\times\to\Gamma
\]
be a valuation. Define
\[
V=\{0\}\cup\{x\in K^\times:v(x)\ge0\},
\]
and
\[
\mathfrak m
=
\{0\}\cup\{x\in K^\times:v(x)>0\}.
\]

\textbf{(a)} Prove that \(V\) is a subring of \(K\).

\textbf{(b)} Prove that \(V\) is a valuation ring.

\textbf{(c)} Prove that \(\mathfrak m\) is the unique maximal ideal of \(V\).

\textbf{(d)} Show that the units of \(V\) are exactly the elements of
valuation \(0\).
\end{problem}

\begin{solution}
If \(x,y\in V\), then
\[
v(xy)=v(x)+v(y)\ge0,
\]
so \(xy\in V\). Also
\[
v(x+y)\ge\min\{v(x),v(y)\}\ge0
\]
whenever \(x+y\ne0\), so \(x+y\in V\). Clearly \(0,1\in V\), hence \(V\) is
a subring.

For any \(x\in K^\times\), either
\[
v(x)\ge0
\]
or
\[
v(x)<0.
\]
In the second case,
\[
v(x^{-1})=-v(x)>0,
\]
so \(x^{-1}\in V\). Therefore \(V\) is a valuation ring.

Now \(x\in V\) is a unit exactly when \(x^{-1}\in V\). This occurs exactly
when
\[
v(x)\ge0
\quad\text{and}\quad
-v(x)\ge0,
\]
hence exactly when
\[
v(x)=0.
\]
Thus the nonunits are precisely the elements of positive valuation together
with \(0\), namely \(\mathfrak m\). Since the nonunits of a local ring form
its unique maximal ideal,
\[
\mathfrak m
\]
is the unique maximal ideal of \(V\).
\end{solution}


\begin{problem}[CP-V-0037]
\label{prob:cp-v-0037}
Let \(V\subseteq K\) be a valuation ring with value group \(\Gamma\).

\textbf{(a)} Prove that every finitely generated ideal of \(V\) is principal.

\textbf{(b)} Prove that \(V\) is Noetherian if and only if either \(V=K\) or
\[
\Gamma\cong\mathbb Z
\]
as an ordered abelian group.

\textbf{(c)} In the nonfield case, deduce that a Noetherian valuation ring is
a discrete valuation ring.
\end{problem}

\begin{solution}
For \textbf{(a)}, let
\[
I=(x_1,\dots,x_n).
\]
Choose \(x_j\) with minimal valuation among the \(v(x_i)\). Then for every
\(i\),
\[
v(x_i/x_j)\ge0,
\]
so
\[
x_i/x_j\in V.
\]
Hence \(x_i\in(x_j)\), and therefore
\[
I=(x_j).
\]

Assume now that \(V\) is Noetherian and not a field. Its maximal ideal
\(\mathfrak m\) is finitely generated, hence principal:
\[
\mathfrak m=(\pi).
\]
Let
\[
\gamma=v(\pi)>0.
\]
If \(\delta>0\) lies in \(\Gamma\), choose \(x\in\mathfrak m\) with
\(v(x)=\delta\). Since \(x\in(\pi)\),
\[
\delta\ge\gamma.
\]
Thus \(\gamma\) is the least positive element of \(\Gamma\).

Now let \(\delta\in\Gamma\) be positive. By the Archimedean argument forced
by principal generation of powers of \(\mathfrak m\), there is a unique
integer \(n\ge0\) with
\[
n\gamma\le\delta<(n+1)\gamma.
\]
Then
\[
0\le\delta-n\gamma<\gamma.
\]
Minimality of \(\gamma\) implies
\[
\delta-n\gamma=0.
\]
Hence every positive element is an integral multiple of \(\gamma\), and
therefore
\[
\Gamma\cong\mathbb Z.
\]

Conversely, suppose
\[
\Gamma\cong\mathbb Z.
\]
Every nonzero ideal \(I\subseteq V\) has a set of valuations
\[
\{v(x):0\ne x\in I\}\subseteq\mathbb Z_{\ge0},
\]
which has a least element \(n\). Choose \(x\in I\) with \(v(x)=n\). For every
\(y\in I\),
\[
v(y/x)\ge0,
\]
so \(y/x\in V\), hence \(y\in(x)\). Therefore
\[
I=(x).
\]
Thus every ideal is principal, so \(V\) is Noetherian.

Therefore a nonfield valuation ring is Noetherian exactly when its value
group is \(\mathbb Z\); this is precisely a discrete valuation ring.
\end{solution}


\begin{problem}[CP-V-0038]
\label{prob:cp-v-0038}
Let \(V\) be a valuation ring of a field \(K\), and let
\[
U=V^\times.
\]
Set
\[
\Gamma=K^\times/U.
\]

\textbf{(a)} Define
\[
\bar x\ge \bar y
\quad\Longleftrightarrow\quad
xy^{-1}\in V.
\]
Prove that this is a well-defined total order on \(\Gamma\) compatible with
the group law.

\textbf{(b)} Let
\[
v:K^\times\to\Gamma
\]
be the quotient map. Prove that
\[
v(x+y)\ge \min\{v(x),v(y)\}
\]
whenever \(x+y\ne0\).

\textbf{(c)} Explain how the valuation ring can be recovered from \(v\).
\end{problem}

\begin{solution}
Suppose \(x'=xu\) and \(y'=yv\) with \(u,v\in U\). Then
\[
x'(y')^{-1}
=
xy^{-1}\,uv^{-1}.
\]
Since \(uv^{-1}\) is a unit of \(V\), multiplication by it does not change
membership in \(V\). Thus the relation is independent of representatives.

Reflexivity follows from \(1\in V\). If
\[
xy^{-1}\in V
\quad\text{and}\quad
yx^{-1}\in V,
\]
then \(xy^{-1}\) is a unit, so \(xU=yU\); hence the order is antisymmetric.
Transitivity follows from
\[
xz^{-1}=(xy^{-1})(yz^{-1})\in V.
\]
Totality is exactly the valuation-ring property: for every
\(xy^{-1}\in K^\times\), either \(xy^{-1}\in V\) or \(yx^{-1}\in V\).

Compatibility with the group law is immediate because
\[
(xz)(yz)^{-1}=xy^{-1}.
\]

Now suppose
\[
v(x)\ge v(y).
\]
Then
\[
xy^{-1}\in V,
\]
so
\[
(x+y)y^{-1}=xy^{-1}+1\in V.
\]
Hence
\[
v(x+y)\ge v(y).
\]
The other case is symmetric, proving
\[
v(x+y)\ge\min\{v(x),v(y)\}.
\]

Finally,
\[
V
=
\{0\}\cup\{x\in K^\times:v(x)\ge0\},
\]
where \(0\in\Gamma\) denotes the class of the units. Thus the ordered quotient
group and quotient map recover the valuation ring.
\end{solution}


\begin{problem}[CP-V-0039]
\label{prob:cp-v-0039}
Let \(k\) be a field and let \(\Gamma\) be a totally ordered abelian group.
Consider the group ring
\[
k[\Gamma]
=
\left\{
\sum_{i=1}^r a_i z^{\gamma_i}:
a_i\in k,\ \gamma_i\in\Gamma
\right\},
\]
with multiplication
\[
z^\gamma z^\delta=z^{\gamma+\delta}.
\]

\textbf{(a)} Prove that \(k[\Gamma]\) is an integral domain.

\textbf{(b)} For nonzero
\[
f=\sum_i a_i z^{\gamma_i},
\]
define
\[
v_0(f)=\min_i\gamma_i.
\]
Prove that
\[
v_0(fg)=v_0(f)+v_0(g)
\]
and
\[
v_0(f+g)\ge\min\{v_0(f),v_0(g)\}
\]
whenever \(f+g\ne0\).

\textbf{(c)} Extend \(v_0\) to the fraction field
\[
K=\operatorname{Frac}(k[\Gamma])
\]
and prove that the resulting valuation has value group exactly \(\Gamma\).
\end{problem}

\begin{solution}
Let
\[
f=\sum_i a_i z^{\gamma_i},
\qquad
g=\sum_j b_j z^{\delta_j},
\]
and put
\[
\gamma_0=\min_i\gamma_i,
\qquad
\delta_0=\min_j\delta_j.
\]
Every exponent appearing in \(fg\) is at least
\[
\gamma_0+\delta_0.
\]
The coefficient of the lowest exponent is nonzero, so
\[
fg\ne0.
\]
Hence \(k[\Gamma]\) is a domain.

The same argument shows that the smallest exponent occurring in \(fg\) is
exactly
\[
\gamma_0+\delta_0,
\]
so
\[
v_0(fg)=v_0(f)+v_0(g).
\]
No exponent smaller than
\[
\min\{v_0(f),v_0(g)\}
\]
can appear in \(f+g\), therefore
\[
v_0(f+g)\ge\min\{v_0(f),v_0(g)\}.
\]

For
\[
x=\frac fg\in K^\times,
\]
define
\[
v(x)=v_0(f)-v_0(g).
\]
If \(f/g=f'/g'\), then
\[
fg'=f'g,
\]
and multiplicativity of \(v_0\) gives
\[
v_0(f)-v_0(g)
=
v_0(f')-v_0(g').
\]
Thus \(v\) is well defined. The multiplicative and ultrametric identities
follow directly from those of \(v_0\).

Finally,
\[
v(z^\gamma)=\gamma
\]
for every \(\gamma\in\Gamma\), so the image of \(v\) is all of \(\Gamma\).
Thus the value group is exactly \(\Gamma\).
\end{solution}


\begin{problem}[CP-V-0040]
\label{prob:cp-v-0040}
Let \(K\) carry a valuation
\[
v:K^\times\to\mathbb R,
\]
and suppose
\[
1+a+b=0
\]
with \(a,b\in K^\times\). Put
\[
\alpha=v(a),
\qquad
\beta=v(b).
\]

\textbf{(a)} Prove that the minimum of
\[
0,\alpha,\beta
\]
must be attained at least twice.

\textbf{(b)} Deduce that
\[
(\alpha,\beta)
\]
lies on the union of the three rays
\[
\{(\lambda,\lambda):\lambda\le0\}
\cup
\{(0,\mu):\mu\ge0\}
\cup
\{(\mu,0):\mu\ge0\}.
\]

\textbf{(c)} In the group-ring valued fields constructed above, show that
every point on these rays is attained, apart from the possible exceptional
vertex phenomenon in residue characteristic \(2\).
\end{problem}

\begin{solution}
Suppose, for example, that
\[
\alpha<\min\{0,\beta\}.
\]
From
\[
a=-(1+b)
\]
we obtain
\[
v(a)=v(1+b)\ge\min\{0,\beta\},
\]
contradicting the strict inequality. The same argument applies if \(0\) or
\(\beta\) is uniquely minimal. Therefore the minimum of
\[
0,\alpha,\beta
\]
is attained at least twice.

There are then exactly three possibilities:
\[
\alpha=\beta\le0,
\qquad
\alpha=0\le\beta,
\qquad
\beta=0\le\alpha.
\]
Hence the image lies on the stated three rays.

Conversely, in a valued field containing monomials \(z^\lambda\) with
\[
v(z^\lambda)=\lambda,
\]
take \(\mu\ge0\) and set
\[
b=z^\mu,
\qquad
a=-1-z^\mu.
\]
Then
\[
1+a+b=0,
\qquad
v(a)=0,
\qquad
v(b)=\mu,
\]
so \((0,\mu)\) occurs. By symmetry \((\mu,0)\) occurs.

For \(\lambda<0\), set
\[
a=z^\lambda,
\qquad
b=-1-z^\lambda.
\]
Then
\[
v(a)=v(b)=\lambda,
\]
so \((\lambda,\lambda)\) occurs.

The origin \((0,0)\) is attained whenever the residue field contains
nonzero \(c,d\) satisfying
\[
1+c+d=0.
\]
Over \(\mathbb F_2\), this can fail, producing the familiar missing-vertex
exception.
\end{solution}



\section*{V/22 --- Chain Complexes}

\begin{problem}[CP-V-0046]
\label{prob:cp-v-0046}
Let
\[
0\longrightarrow C'_\bullet
\xrightarrow{\alpha}
C_\bullet
\xrightarrow{\beta}
C''_\bullet
\longrightarrow0
\]
be a short exact sequence of chain complexes.

\textbf{(a)} Construct the connecting homomorphism
\[
\partial_n:H_n(C''_\bullet)\longrightarrow H_{n-1}(C'_\bullet).
\]

\textbf{(b)} Prove that these maps fit into a long exact sequence
\[
\cdots\to
H_n(C'_\bullet)\to H_n(C_\bullet)\to H_n(C''_\bullet)
\xrightarrow{\partial_n}
H_{n-1}(C'_\bullet)\to\cdots.
\]

\textbf{(c)} Explain why a short exact sequence
\[
0\to N'\to N\to N''\to0
\]
and a free resolution \(F_\bullet\to M\) yield the long exact Tor sequence
\[
\cdots\to
\operatorname{Tor}_n^A(M,N')
\to
\operatorname{Tor}_n^A(M,N)
\to
\operatorname{Tor}_n^A(M,N'')
\to
\operatorname{Tor}_{n-1}^A(M,N')
\to\cdots.
\]
\end{problem}

\begin{solution}
Let
\[
[z'']\in H_n(C''_\bullet)
\]
with \(d z''=0\). Choose \(z\in C_n\) with
\[
\beta(z)=z''.
\]
Then
\[
\beta(dz)=d\beta(z)=dz''=0,
\]
so by exactness there is \(z'\in C'_{n-1}\) with
\[
\alpha(z')=dz.
\]
Moreover,
\[
\alpha(dz')=d\alpha(z')=d^2z=0.
\]
Since \(\alpha\) is injective,
\[
dz'=0.
\]
Define
\[
\partial_n([z''])=[z'].
\]
Changing the lift \(z\), or changing \(z''\) by a boundary, changes \(z'\)
only by a boundary, so the map is well defined.

Exactness at each term follows by the standard diagram chase. For example,
an element of \(H_n(C''_\bullet)\) lies in \(\ker\partial_n\) precisely when
one can modify its lift in \(C_n\) to make it a cycle, which is precisely
when it comes from \(H_n(C_\bullet)\). The other terms are analogous.
Therefore the homology groups form a long exact sequence.

For \textbf{(c)}, since every \(F_i\) is free, hence flat, tensoring
\[
0\to N'\to N\to N''\to0
\]
with \(F_i\) remains exact:
\[
0\to F_i\otimes_A N'
\to F_i\otimes_A N
\to F_i\otimes_A N''
\to0.
\]
Degree by degree this gives a short exact sequence of chain complexes
\[
0\to F_\bullet\otimes_A N'
\to F_\bullet\otimes_A N
\to F_\bullet\otimes_A N''
\to0.
\]
Taking the long exact homology sequence and using
\[
H_i(F_\bullet\otimes_A N)=\operatorname{Tor}_i^A(M,N)
\]
gives the required long exact Tor sequence.
\end{solution}


\begin{problem}[CP-V-0110]
\label{prob:cp-v-0110}
Let \(f,g:C_\bullet\to D_\bullet\) be chain maps. Suppose there are homomorphisms
\[
h_n:C_n\to D_{n+1}
\]
such that
\[
f-g=d_Dh+hd_C.
\]
Prove that \(f\) and \(g\) induce the same map on homology.
\end{problem}

\begin{solution}
Let \(z\in C_n\) be a cycle. Then
\[
(f-g)(z)=d_Dh(z)+h(d_Cz)=d_Dh(z),
\]
because \(d_Cz=0\). Thus \(f(z)-g(z)\) is a boundary in \(D_n\). Hence
\[
[f(z)]=[g(z)]
\]
in \(H_n(D_\bullet)\). Therefore the induced maps on homology agree in every degree.
\end{solution}



\section*{V/23 --- Free Resolutions}

\begin{problem}[CP-V-0041]
\label{prob:cp-v-0041}
Let
\[
A=k[x,y]/(xy),
\qquad
\mathfrak m=(x,y),
\qquad
k=A/\mathfrak m.
\]

\textbf{(a)} Construct a free resolution of \(k\) by finitely generated free
\(A\)-modules.

\textbf{(b)} Compute
\[
\operatorname{Ext}_A^i(k,A)
\]
for all \(i\ge0\).

\textbf{(c)} Compute
\[
\operatorname{Tor}_i^A(k,k)
\]
for all \(i\ge0\).
\end{problem}

\begin{solution}
The natural surjection
\[
A\twoheadrightarrow k
\]
has kernel \(\mathfrak m=(x,y)\). Hence begin with
\[
A^2\xrightarrow{d_1}A\to k\to0,
\qquad
d_1(a,b)=ax+by.
\]
Thus
\[
d_1=\begin{bmatrix}x&y\end{bmatrix}.
\]

In \(A\),
\[
\operatorname{Ann}(x)=(y),
\qquad
\operatorname{Ann}(y)=(x).
\]
Therefore
\[
\ker d_1=(y)\times(x),
\]
which is the image of
\[
d_2=
\begin{bmatrix}
y&0\\
0&x
\end{bmatrix}.
\]
Similarly,
\[
\ker d_2=(x)\times(y),
\]
which is the image of
\[
d_3=
\begin{bmatrix}
x&0\\
0&y
\end{bmatrix}.
\]
Repeating gives a \(2\)-periodic resolution
\[
\cdots\to A^2
\xrightarrow{\left[\begin{smallmatrix}y&0\\0&x\end{smallmatrix}\right]}
A^2
\xrightarrow{\left[\begin{smallmatrix}x&0\\0&y\end{smallmatrix}\right]}
A^2
\xrightarrow{\left[\begin{smallmatrix}y&0\\0&x\end{smallmatrix}\right]}
A^2
\xrightarrow{\left[\begin{smallmatrix}x&y\end{smallmatrix}\right]}
A
\to k\to0.
\]

Applying \(\operatorname{Hom}_A(-,A)\) gives
\[
0\to A\xrightarrow{\delta^0}A^2\xrightarrow{\delta^1}A^2
\xrightarrow{\delta^2}\cdots,
\]
where
\[
\delta^0(a)=(xa,ya),
\]
and the later maps alternate between the two diagonal matrices above.

Now
\[
\ker\delta^0=(x)\cap(y)=0,
\]
so
\[
\operatorname{Ext}_A^0(k,A)=0.
\]
Also
\[
\ker\delta^1=(x)\times(y),
\]
while
\[
\operatorname{im}\delta^0=\{(xa,ya):a\in A\}.
\]
The quotient is one-dimensional over \(k\), hence
\[
\operatorname{Ext}_A^1(k,A)\cong k.
\]
For \(i\ge2\), kernels and images alternate and agree, so
\[
\operatorname{Ext}_A^i(k,A)=0.
\]

Thus
\[
\operatorname{Ext}_A^i(k,A)\cong
\begin{cases}
0,&i=0,\\
k,&i=1,\\
0,&i\ge2.
\end{cases}
\]

Finally, tensor the free resolution with \(k\). Since \(x\) and \(y\) act
trivially on \(k\), every differential becomes zero. Therefore
\[
\operatorname{Tor}_0^A(k,k)\cong k,
\qquad
\operatorname{Tor}_i^A(k,k)\cong k^2
\quad(i\ge1).
\]
\end{solution}


\begin{problem}[CP-V-0077]
\label{prob:cp-v-0077}
Find a free resolution of \(\mathbb Z/12\mathbb Z\) as a \(\mathbb Z\)-module.
\end{problem}

\begin{solution}
The standard resolution is
\[
0 \longrightarrow \mathbb Z \xrightarrow{\cdot 12} \mathbb Z \longrightarrow \mathbb Z/12\mathbb Z \longrightarrow 0.
\]
The quotient map has kernel \(12\mathbb Z\), which is the image of multiplication by \(12\).
\end{solution}


\begin{problem}[CP-V-0078]
\label{prob:cp-v-0078}
Find a free resolution of \(k[x]/(x^3)\) as a \(k[x]\)-module.
\end{problem}

\begin{solution}
We have
\[
0 \longrightarrow k[x] \xrightarrow{\cdot x^3} k[x] \longrightarrow k[x]/(x^3) \longrightarrow 0.
\]
This is exact because the kernel of the quotient map is the principal ideal \((x^3)\), which is isomorphic to \(k[x]\).
\end{solution}


\begin{problem}[CP-V-0079]
\label{prob:cp-v-0079}
Describe a free resolution of \(k\) as a \(k[x,y]\)-module.
\end{problem}

\begin{solution}
Since
\[
k \cong k[x,y]/(x,y),
\]
a free resolution is
\[
0 \longrightarrow A \xrightarrow{\begin{pmatrix}-y\\x\end{pmatrix}} A^2
\xrightarrow{(x\;\;y)} A \longrightarrow k \longrightarrow 0,
\qquad A=k[x,y].
\]
The two relations are \(x=0\) and \(y=0\), and the relation among them is \((-y,x)\).
\end{solution}


\begin{problem}[CP-V-0082]
\label{prob:cp-v-0082}
Explain why \(k[x]/(x)\) has a trivial free resolution over \(k[x]/(x)\), but a nontrivial free resolution over \(k[x]\).
\end{problem}

\begin{solution}
As a module over \(k[x]/(x)\cong k\), it is simply \(k\), which is free of rank \(1\), so the resolution is trivial.
As a module over \(k[x]\), it is \(k[x]/(x)\), which is not free over \(k[x]\), but it has the nontrivial free resolution
\[
0 \longrightarrow k[x] \xrightarrow{\cdot x} k[x] \longrightarrow k[x]/(x) \longrightarrow 0.
\]
\end{solution}



\section*{V/24 --- Syzygies}

\begin{problem}[CP-V-0111]
\label{prob:cp-v-0111}
Let
\[
0\to K\to F\to M\to0
\]
be exact with \(F\) free. The module \(K\) is called a first syzygy of \(M\). If
\[
0\to K'\to F'\to M\to0
\]
is another such sequence with \(F'\) free, prove Schanuel's lemma:
\[
K\oplus F'\cong K'\oplus F.
\]
\end{problem}

\begin{solution}
Form the pullback
\[
P=\{(x,y)\in F\oplus F':\bar x=\bar y\in M\}.
\]
Projection \(P\to F\) has kernel naturally isomorphic to \(K'\), and since \(F\) is free the exact sequence
\[
0\to K'\to P\to F\to0
\]
splits. Hence
\[
P\cong K'\oplus F.
\]
Similarly, projection \(P\to F'\) has kernel \(K\), and because \(F'\) is free,
\[
P\cong K\oplus F'.
\]
Combining the two splittings gives
\[
K\oplus F'\cong K'\oplus F.
\]
\end{solution}


\begin{problem}[CP-V-0112]
\label{prob:cp-v-0112}
Let \(A=k[x,y]\) and \(I=(x^2,xy,y^2)\). Find generators for the first syzygy module of the generating set \(x^2,xy,y^2\).
\end{problem}

\begin{solution}
Consider
\[
A^3\longrightarrow I,\qquad (a,b,c)\longmapsto ax^2+bxy+cy^2.
\]
Two evident syzygies are
\[
(-y,x,0),\qquad(0,-y,x).
\]
They generate all syzygies. Indeed, if
\[
ax^2+bxy+cy^2=0,
\]
then setting \(y=0\) shows \(a\) is divisible by \(y\), say \(a=-yu\). Setting \(x=0\) shows \(c\) is divisible by \(x\), say \(c=xv\). Substitution gives
\[
xy(-ux+b+vy)=0,
\]
so
\[
b=ux-vy.
\]
Therefore
\[
(a,b,c)=u(-y,x,0)+v(0,-y,x).
\]
\end{solution}



\section*{V/25 --- Minimal Resolutions}

\begin{problem}[CP-V-0042]
\label{prob:cp-v-0042}
Let \((A,\mathfrak m)\) be a Noetherian local ring, let
\[
k=A/\mathfrak m,
\]
and let \(M\) be a finitely generated \(A\)-module.

A free resolution
\[
\cdots\longrightarrow F_2\xrightarrow{d_2}F_1\xrightarrow{d_1}F_0
\longrightarrow M\longrightarrow0
\]
by finitely generated free modules is called \emph{minimal} if
\[
d_i\otimes_A k=0
\]
for every \(i\ge1\).

Prove that every finitely generated \(A\)-module admits a minimal free
resolution.
\end{problem}

\begin{solution}
Choose a minimal free cover
\[
\varepsilon_0:F_0\twoheadrightarrow M
\]
as in CP-V-0021. Then
\[
K_1:=\ker(\varepsilon_0)
\subseteq
\mathfrak mF_0.
\]
Since \(A\) is Noetherian and \(F_0\) is finitely generated, \(K_1\) is
finitely generated.

Choose a minimal free cover
\[
\varepsilon_1:F_1\twoheadrightarrow K_1.
\]
Composing with the inclusion \(K_1\hookrightarrow F_0\) gives
\[
d_1:F_1\to F_0
\]
with
\[
\operatorname{im}d_1=K_1
\]
and
\[
\operatorname{im}d_1\subseteq\mathfrak mF_0.
\]
Therefore
\[
d_1\otimes_Ak=0.
\]

Now put
\[
K_2:=\ker d_1.
\]
Again \(K_2\) is finitely generated because it is a submodule of the
finitely generated free module \(F_1\). Choose a minimal free cover
\[
\varepsilon_2:F_2\twoheadrightarrow K_2
\]
and compose with \(K_2\hookrightarrow F_1\) to obtain
\[
d_2:F_2\to F_1.
\]
Minimality gives
\[
\operatorname{im}d_2\subseteq\mathfrak mF_1,
\]
hence
\[
d_2\otimes_Ak=0.
\]

Proceed inductively. At stage \(n\), define
\[
K_{n+1}=\ker d_n.
\]
Noetherianity ensures that \(K_{n+1}\) is finitely generated. Choose a
minimal free cover
\[
F_{n+1}\twoheadrightarrow K_{n+1},
\]
and compose with the inclusion \(K_{n+1}\hookrightarrow F_n\). This produces
an exact free resolution in which
\[
\operatorname{im}d_{n+1}\subseteq\mathfrak mF_n
\]
for every \(n\), equivalently
\[
d_{n+1}\otimes_Ak=0.
\]

Thus minimal free resolutions exist.
\end{solution}


\begin{problem}[CP-V-0043]
\label{prob:cp-v-0043}
Let \((A,\mathfrak m)\) be a Noetherian local ring, let
\[
k=A/\mathfrak m,
\]
and let \(M\) be a finitely generated \(A\)-module with finite projective
dimension.

Prove that:

\textbf{(a)} every minimal free resolution of \(M\) has length
\[
\operatorname{pd}_A M;
\]

\textbf{(b)}
\[
\operatorname{Tor}_i^A(k,M)\cong F_i\otimes_Ak
\]
for a minimal free resolution \(F_\bullet\to M\);

\textbf{(c)}
\[
\operatorname{pd}_A M
=
\min\left\{
i\ge0:
\operatorname{Tor}_{i+1}^A(k,M)=0
\right\}.
\]
\end{problem}

\begin{solution}
Let
\[
\cdots\to F_2\to F_1\to F_0\to M\to0
\]
be a minimal free resolution. By definition,
\[
d_i\otimes_Ak=0
\]
for all \(i\). Hence after tensoring with \(k\), the entire complex has zero
differentials. Therefore
\[
\operatorname{Tor}_i^A(k,M)
\cong
F_i\otimes_Ak.
\]
Since \(F_i\) is finitely generated free,
\[
F_i\otimes_Ak=0
\iff
F_i=0.
\]

Let
\[
n=\operatorname{pd}_A M.
\]
There exists some free resolution of length \(n\), so
\[
\operatorname{Tor}_i^A(k,M)=0
\qquad(i>n).
\]
Using the minimal resolution,
\[
F_i\otimes_Ak=0
\qquad(i>n),
\]
hence
\[
F_i=0
\qquad(i>n).
\]
Thus the minimal resolution has length at most \(n\).

If \(F_n=0\), then the minimal resolution would already terminate in degree
\(n-1\), producing a free resolution of length at most \(n-1\), contrary to
the definition of \(n\). Hence
\[
F_n\ne0.
\]
So every minimal free resolution has length exactly \(n\).

Consequently,
\[
\operatorname{Tor}_n^A(k,M)\ne0,
\qquad
\operatorname{Tor}_{n+1}^A(k,M)=0.
\]
Therefore
\[
\operatorname{pd}_A M
=
\min\left\{
i\ge0:
\operatorname{Tor}_{i+1}^A(k,M)=0
\right\}.
\]
\end{solution}



\section*{V/26 --- The Tor Functor}

\begin{problem}[CP-V-0047]
\label{prob:cp-v-0047}
Let \(n,m\ge1\).

\textbf{(a)} Prove that for every abelian group \(N\),
\[
\operatorname{Tor}_1^{\mathbb Z}(\mathbb Z/n\mathbb Z,N)
\cong
N[n]
:=
\{x\in N:nx=0\}.
\]

\textbf{(b)} Deduce that
\[
\operatorname{Tor}_1^{\mathbb Z}
(\mathbb Z/n\mathbb Z,\mathbb Z/m\mathbb Z)
\cong
\mathbb Z/\gcd(n,m)\mathbb Z.
\]

\textbf{(c)} Compare this with
\[
(\mathbb Z/n\mathbb Z)\otimes_{\mathbb Z}
(\mathbb Z/m\mathbb Z).
\]
\end{problem}

\begin{solution}
Use the free resolution
\[
0\longrightarrow\mathbb Z
\xrightarrow{\cdot n}
\mathbb Z
\longrightarrow
\mathbb Z/n\mathbb Z
\longrightarrow0.
\]
Tensoring with \(N\) gives the complex
\[
0\longrightarrow N
\xrightarrow{\cdot n}
N
\longrightarrow0.
\]
Therefore
\[
\operatorname{Tor}_1^{\mathbb Z}(\mathbb Z/n\mathbb Z,N)
=
\ker(N\xrightarrow{\cdot n}N)
=
N[n].
\]

For \(N=\mathbb Z/m\mathbb Z\), the kernel of multiplication by \(n\) has
exactly
\[
d=\gcd(n,m)
\]
elements and is cyclic. Thus
\[
\operatorname{Tor}_1^{\mathbb Z}
(\mathbb Z/n,\mathbb Z/m)
\cong
\mathbb Z/d\mathbb Z.
\]

Similarly,
\[
(\mathbb Z/n)\otimes_{\mathbb Z}(\mathbb Z/m)
\cong
(\mathbb Z/m)/n(\mathbb Z/m)
\cong
\mathbb Z/d\mathbb Z.
\]
Thus in this particular pair of cyclic groups, degree-zero tensor product
and degree-one Tor happen to be isomorphic, although they arise from
different parts of the tensor complex.
\end{solution}


\begin{problem}[CP-V-0048]
\label{prob:cp-v-0048}
Let \(A\) be a ring and let \(I,J\subseteq A\) be ideals. Prove the natural
isomorphism
\[
\operatorname{Tor}_1^A(A/I,A/J)
\cong
\frac{I\cap J}{IJ}.
\]

Use it to compute:

\textbf{(a)}
\[
\operatorname{Tor}_1^{k[x]}
\bigl(k[x]/(x),k[x]/(x)\bigr);
\]

\textbf{(b)}
\[
\operatorname{Tor}_1^{k[x]}
\bigl(k[x]/(x),k[x]/(x-1)\bigr).
\]
\end{problem}

\begin{solution}
Start with
\[
0\longrightarrow I\longrightarrow A\longrightarrow A/I\longrightarrow0.
\]
Tensor with \(A/J\). The long exact Tor sequence begins
\[
0\longrightarrow
\operatorname{Tor}_1^A(A/I,A/J)
\longrightarrow
I\otimes_A A/J
\longrightarrow
A\otimes_A A/J.
\]
Using
\[
I\otimes_A A/J\cong I/IJ
\]
and
\[
A\otimes_A A/J\cong A/J,
\]
the relevant map is
\[
I/IJ\longrightarrow A/J,
\qquad
i+IJ\longmapsto i+J.
\]
Its kernel consists precisely of those \(i\in I\cap J\), modulo \(IJ\).
Hence
\[
\operatorname{Tor}_1^A(A/I,A/J)
\cong
\frac{I\cap J}{IJ}.
\]

For \textbf{(a)}, take
\[
I=J=(x).
\]
Then
\[
I\cap J=(x),
\qquad
IJ=(x^2),
\]
so
\[
\operatorname{Tor}_1^{k[x]}(k,k)
\cong
(x)/(x^2)
\cong k.
\]

For \textbf{(b)}, the ideals \((x)\) and \((x-1)\) are comaximal, hence
\[
(x)\cap(x-1)=(x)(x-1).
\]
Therefore
\[
\operatorname{Tor}_1^{k[x]}
\bigl(k[x]/(x),k[x]/(x-1)\bigr)
=0.
\]
\end{solution}



\section*{V/27 --- The Ext Functor}

\begin{problem}[CP-V-0045]
\label{prob:cp-v-0045}
Let \((A,\mathfrak m)\) be a local ring with residue field
\[
k=A/\mathfrak m,
\]
and suppose
\[
0\longrightarrow K\xrightarrow{\,f\,}F\xrightarrow{\,g\,}M\longrightarrow0
\]
is exact, where \(K\) and \(F\) are finite free \(A\)-modules and \(g\) is a
minimal free cover.

Prove that for every \(i\ge0\), the induced homomorphism
\[
f_*:\operatorname{Ext}_A^i(k,K)\longrightarrow
\operatorname{Ext}_A^i(k,F)
\]
is zero.
\end{problem}

\begin{solution}
Because \(g\) is minimal,
\[
\ker g\subseteq\mathfrak mF.
\]
Since
\[
\operatorname{im}f=\ker g,
\]
we have
\[
\operatorname{im}f\subseteq\mathfrak mF.
\]

Choose bases
\[
K\cong A^n,\qquad F\cong A^m.
\]
Then \(f\) is represented by a matrix
\[
C=(c_{ij})
\]
whose entries all lie in \(\mathfrak m\).

For every \(i\ge0\),
\[
\mathfrak m\operatorname{Ext}_A^i(k,A)=0.
\]
Indeed, the \(A\)-module action on \(\operatorname{Ext}_A^i(k,A)\) may be
read from the first variable \(k\); every element of \(\mathfrak m\) acts
trivially on
\[
k=A/\mathfrak m.
\]
Hence it acts trivially on the Ext group.

Using
\[
\operatorname{Ext}_A^i(k,K)
\cong
\operatorname{Ext}_A^i(k,A)^n
\]
and
\[
\operatorname{Ext}_A^i(k,F)
\cong
\operatorname{Ext}_A^i(k,A)^m,
\]
the map \(f_*\) is represented by the same matrix \(C\). Every entry
\(c_{ij}\in\mathfrak m\) acts as zero, so
\[
f_*=0
\]
for all \(i\ge0\).
\end{solution}


\begin{problem}[CP-V-0113]
\label{prob:cp-v-0113}
Compute
\[
\operatorname{Ext}^1_{\mathbb Z}(\mathbb Z/n\mathbb Z,\mathbb Z/m\mathbb Z).
\]
\end{problem}

\begin{solution}
Use the free resolution
\[
0\to\mathbb Z\xrightarrow{\cdot n}\mathbb Z\to\mathbb Z/n\mathbb Z\to0.
\]
Applying
\[
\operatorname{Hom}_{\mathbb Z}(-,\mathbb Z/m\mathbb Z)
\]
gives
\[
\mathbb Z/m\mathbb Z\xrightarrow{\cdot n}\mathbb Z/m\mathbb Z.
\]
Hence
\[
\operatorname{Ext}^1_{\mathbb Z}(\mathbb Z/n,\mathbb Z/m)
\cong
\frac{\mathbb Z/m}{n(\mathbb Z/m)}
\cong
\mathbb Z/\gcd(n,m)\mathbb Z.
\]
\end{solution}



\section*{V/28 --- Derived-Functor Viewpoint}

\begin{problem}[CP-V-0044]
\label{prob:cp-v-0044}
Let \(A\) be a Noetherian local ring and let \(M\ne0\) be a finitely generated
\(A\)-module with finite projective dimension. Prove the
Auslander--Buchsbaum formula
\[
\operatorname{pd}_A M+\operatorname{depth}(M)
=
\operatorname{depth}(A).
\]
You may use the characterization
\[
\operatorname{depth}(N)
=
\min\left\{
i\ge0:
\operatorname{Ext}_A^i(k,N)\ne0
\right\},
\qquad
k=A/\mathfrak m.
\]
\end{problem}

\begin{solution}
Set
\[
s=\operatorname{depth}(A),
\qquad
n=\operatorname{pd}_A M.
\]
We argue by induction on \(n\).

If \(n=0\), then \(M\) is projective. Since \(A\) is local and \(M\) is
finitely generated, \(M\) is free:
\[
M\cong A^r.
\]
Hence
\[
\operatorname{Ext}_A^i(k,M)
\cong
\bigl(\operatorname{Ext}_A^i(k,A)\bigr)^r,
\]
so
\[
\operatorname{depth}(M)=s.
\]

Now assume \(n>0\). Take the beginning of a minimal free resolution:
\[
0\longrightarrow K\longrightarrow F\longrightarrow M\longrightarrow0,
\]
with \(F\) finitely generated free and
\[
\operatorname{pd}_A K=n-1.
\]
By the induction hypothesis,
\[
\operatorname{depth}(K)
=
s-(n-1)
=
s-n+1.
\]
Put
\[
t=\operatorname{depth}(K)=s-n+1.
\]

Because the presentation is minimal, the map
\[
K\to F
\]
has image in \(\mathfrak mF\). The induced maps
\[
\operatorname{Ext}_A^i(k,K)
\longrightarrow
\operatorname{Ext}_A^i(k,F)
\]
are therefore zero. Hence the long exact Ext sequence breaks into short exact
pieces
\[
0\to
\operatorname{Ext}_A^i(k,F)
\to
\operatorname{Ext}_A^i(k,M)
\to
\operatorname{Ext}_A^{i+1}(k,K)
\to0.
\]

Since \(F\) is free,
\[
\operatorname{depth}(F)=s.
\]
Thus for \(i<t-1\),
\[
\operatorname{Ext}_A^i(k,F)=0,
\qquad
\operatorname{Ext}_A^{i+1}(k,K)=0,
\]
so
\[
\operatorname{Ext}_A^i(k,M)=0.
\]

At \(i=t-1\), the two Ext terms involving \(F\) still vanish, while
\[
\operatorname{Ext}_A^t(k,K)\ne0
\]
by the definition of \(t\). Therefore
\[
\operatorname{Ext}_A^{t-1}(k,M)\ne0.
\]
Hence
\[
\operatorname{depth}(M)=t-1=s-n.
\]
Thus
\[
\operatorname{pd}_A M+\operatorname{depth}(M)
=
n+(s-n)
=
s
=
\operatorname{depth}(A).
\]
\end{solution}


\begin{problem}[CP-V-0114]
\label{prob:cp-v-0114}
Let
\[
0\to A'\to A\to A''\to0
\]
be a short exact sequence of \(R\)-modules. Explain why the long exact sequence of \(\operatorname{Tor}\) measures the failure of tensor product to be left exact, while the long exact sequence of \(\operatorname{Ext}\) measures the failure of \(\operatorname{Hom}\) to be right exact.
\end{problem}

\begin{solution}
For a fixed module \(M\), tensor product
\[
M\otimes_R-
\]
is right exact but may fail to preserve injectivity. Its left derived functors are
\[
\operatorname{Tor}_i^R(M,-),
\]
and the long exact sequence begins
\[
\operatorname{Tor}_1^R(M,A'')\to M\otimes_R A'\to M\otimes_R A\to M\otimes_R A''\to0.
\]
Thus \(\operatorname{Tor}_1\) records precisely the obstruction to left exactness.

Dually,
\[
\operatorname{Hom}_R(M,-)
\]
is left exact but may fail to preserve surjectivity. Its right derived functors are
\[
\operatorname{Ext}_R^i(M,-),
\]
and the long exact sequence contains
\[
\operatorname{Hom}_R(M,A)\to\operatorname{Hom}_R(M,A'')
\to\operatorname{Ext}_R^1(M,A').
\]
Thus \(\operatorname{Ext}^1\) records the obstruction to right exactness. Higher Tor and Ext measure the successive higher failures of exactness.
\end{solution}
